% Fate's Edge – The Gray Wanderer's Grimoire
% A Thematic Guide for Cantors, Invokers, and Runekeepers
% Compile with XeLaTeX or LuaLaTeX

\documentclass[10pt,a5paper]{book}
\usepackage[top=1.5cm,bottom=1.5cm,left=1.2cm,right=1.2cm]{geometry}
\usepackage{fontspec}
\setmainfont{TeX Gyre Pagella}
\usepackage{titlesec}
\usepackage{tcolorbox}
\tcbuselibrary{skins,breakable}
\usepackage{multicol}
\usepackage{hyperref}
\hypersetup{colorlinks=true,linkcolor=blue}
\usepackage{array}
\usepackage{longtable}
\usepackage{epigraph}
\setlength{\epigraphwidth}{0.7\textwidth}
\renewcommand{\epigraphflush}{flushright}
\usepackage{pgfornament}
\usepackage{lettrine}

% Chapter styling
\titleformat{\chapter}[display]
{\Huge\bfseries\scshape}
{\flushright\pgfornament[width=0.2\textwidth]{88}\\\vspace{-1cm}}
{0pt}
{\hfill}
\titlespacing*{\chapter}{0pt}{-20pt}{40pt}

% Section styling
\titleformat{\section}
{\Large\bfseries}
{\thesection}{0.5em}{}
\titlespacing*{\section}{0pt}{12pt}{6pt}

% Subsection
\titleformat{\subsection}
{\normalsize\bfseries\itshape}
{\thesubsection}{0.5em}{}
\titlespacing*{\subsection}{0pt}{8pt}{4pt}

% Custom environments
\newtcolorbox{patronbox}[2][]{
    colback=gray!5,
    colframe=gray!70,
    fonttitle=\bfseries,
    title=⛧ #2 ⛧,
    enhanced,
    attach boxed title to top left={yshift=-2mm,xshift=3mm},
    boxed title style={colback=gray!20},
    #1
}

\newtcolorbox{ritualnote}{
    colback=yellow!5,
    colframe=orange!70,
    before skip=0.5cm,
    after skip=0.5cm
}

\newtcolorbox{warningbox}{
    colback=red!5,
    colframe=red!70,
    before skip=0.5cm,
    after skip=0.5cm
}

% Tags for magic
\newcommand{\ritualtag}[1]{\textsc{#1}}

\begin{document}

\frontmatter
\thispagestyle{empty}
\vspace*{2cm}
\begin{center}
{\Huge\scshape The Gray Wanderer's Grimoire}\\[0.5cm]
{\Large A Thematic Guide for Cantors, Invokers, and Runekeepers}\\[1cm]
\pgfornament[width=0.6\textwidth]{88}\\[1cm]
{\itshape “You want power? Fine. But don’t come crying when the debt comes due.”}
\end{center}
\vfill
\clearpage

\tableofcontents
\clearpage

\mainmatter

\chapter{Introduction – A Voice from the Crossroads}
\epigraph{“I have worn the chains of a dozen Patrons and broken most of them. Twice.”}{\textit{— The Gray Wanderer}}

This grimoire is not a rulebook. It is a \textbf{companion} – a set of lenses through which to see the magical paths of \textit{Fate's Edge}. We will explore:

\begin{itemize}
    \item The \textbf{lore and roleplay} of Cantors, Invokers, and Runekeepers – who they are, how they think, why they risk annihilation.
    \item How to handle \textbf{Obligation} – not as a punishment, but as a story engine. You will find patron‑specific ideas, intrusions, and acts of service.
    \item \textbf{Rituals without a focus} – when your Codex is ash, your Symbol cracked, your voice alone. Components matter again.
\end{itemize}

The voice you hear belongs to no one and everyone – a wanderer who has stood at the crossroads of every Patron’s attention. Call me the \textbf{Gray Wanderer}. These notes are my legacy. Use them, lose them, or burn them for warmth. I won’t mind.

\clearpage

\chapter{The Three Paths – Lore and Roleplay}

\section{The Runekeeper – Scholar of the Inked Vow}
\textit{“I do not beg. I record.”}

Runekeepers are archivists of the divine. Their power comes from knowing the \textit{true name} of a Rite, the correct gesture, the exact syllable, the precise weight of the offering. They carry a \textbf{Codex} – a personal grimoire – and a \textbf{Thiasos} (familiar), a fragment of their Patron’s attention made flesh.

\subsection{Lore – Stone and Bark}
In Aeler, Runekeepers are \textit{Stone‑Scribes}, carving Rites into vault walls with their own blood. In Valewood, they weave Rites into living bark. Every Runekeeper knows that forgetting a single line can kill you.

\subsection{Roleplay Hooks}
\begin{itemize}
    \item Your Codex is incomplete – the last page was torn out. What was on it?
    \item Your Thiasos is a sarcastic rat that only speaks in legal jargon.
    \item A rival Runekeeper claims you stole a Rite from their Codex. They want it back – or your blood.
\end{itemize}

\subsection{Codex \& Thiasos in Play}
\begin{itemize}
    \item \textbf{Losing your Codex:} You cannot access any Rite until you recover it or spend 2 XP to recreate it from memory (requires a day of quiet and a Lore+Spirit test, DV 4).
    \item \textbf{Losing your Thiasos:} You lose the Patron’s Gift (Imbuement) until you perform a ritual of re‑attunement (see Chapter 3).
\end{itemize}

\section{The Invoker – Borrower of Stolen Fire}
\textit{“I don’t serve. I rent.”}

Invokers carry \textbf{Symbols} – physical anchors to Patrons they never fully swore to. They are gamblers, thieves, and desperate souls. In Silkstrand they are called \textit{Bridge‑Runners}. In Acasia, hedge‑witches with a dozen fetishes in their coat.

\subsection{Lore – The Cracked Seal}
Every Invoker knows that a Symbol can break. When it does, the Patron might come to collect in person – or send something smaller, sharper, and hungrier.

\subsection{Roleplay Hooks}
\begin{itemize}
    \item One of your Symbols was stolen by a former student. Now they use it against you.
    \item You have never met your Patron – you found the Symbol in a ruin. Does the Patron even know you exist?
    \item You carry four Symbols (the limit before friction). Every morning, roll a die; on a 1, two Symbols clash. Mark 1 Fatigue.
\end{itemize}

\subsection{Symbols and Obligation}
A \textbf{well‑maintained Symbol} (polished, offered to weekly) reduces Obligation by 1 when you use it (once per session). A \textbf{compromised Symbol} (cracked, dull) can still be used with \textit{Crack the Seal}, but the GM may add a thematic curse (e.g., the Patron’s sigil appears on your skin for a day).

\section{The Cantor – Voice of the Unending Refrain}
\textit{“You think you need a lute? My larynx is older than any tree. Hum, and the world will listen.”}

Cantors sing power into being. They are the most public of the paths – you cannot hide a hymn. Their magic is as natural as breathing, yet as addictive as poppy wine.

\subsection{Lore – The Body Remembers}
Cantors carry no Codex. No Thiasos. Their instrument is their own flesh. A Cantor who loses their voice (sore throat, gag, strangulation) cannot cast Rites – but they can hum through clenched teeth, tap a rhythm on their ribs, or whistle. The body adapts. The body \textit{remembers}.

\subsection{Roleplay Hooks}
\begin{itemize}
    \item You once sang a forbidden note, and now a fragment of a dead Cantor lives in your head.
    \item Your voice is unmistakable – enemies recognise you by your timbre.
    \item You are trying to write a new Song, but every attempt ends in a headache and a page of gibberish. The Song is writing \textit{you}.
\end{itemize}

\subsection{Enchanted Instruments (Optional, Often Cursed)}
A Cantor may acquire an enchanted instrument. It is never bought – it is found, won, or inherited. And it always carries a \textbf{subtle curse}.

\begin{center}
\begin{longtable}{p{2.5cm} p{3.5cm} p{5cm}}
\caption{Sample Enchanted Instruments}\\
\toprule
Instrument & Benefit & Curse \\
\midrule
Silver Flute & +1 die to Performance & Each use attracts a spectral audience; GM gains +1 SB for social complications. \\
Drum of the Deep & +1 Effect on sea shanties & Every third beat summons a minor water spirit (harmless but distracting). \\
Harp of Broken Strings & Re‑roll one failed Performance per session & One string always out of tune; hearers test Spirit+Resolve (DV 3) or become melancholic. \\
Whistle of the Pale Shepherd & Sound carries to Far range regardless of noise & After playing, you forget the next face you see (GM chooses). \\
\bottomrule
\end{longtable}
\end{center}

\subsection{Performance as Position Modifier}
When a Cantor sings to \textit{support} (not to cast a Rite), treat it as a \textbf{Performance Maneuver}:

\begin{tcolorbox}[colback=green!5, colframe=green!50]
\textbf{Performance Maneuver (Action)} \\
\textit{DV 3 (Controlled).} \\
\textbf{Success:} All allies in Near gain +1 Position on their next action. Lasts until Cantor’s next turn. \\
\textbf{Partial:} Same benefit, but Cantor marks 1 Fatigue. \\
\textbf{Miss:} Voice cracks – all allies in Near suffer –1 Position on their next action.
\end{tcolorbox}

\clearpage

\chapter{Obligation – The Debt You Carry}
\epigraph{“Every Rite is a receipt. Every Patron is a creditor. Pay on time, or they’ll repossess your soul.”}{\textit{— The Gray Wanderer}}

Obligation is \textbf{not a punishment}. It is a story engine. When your Obligation exceeds \texttt{Spirit + Presence}, you start marking Fatigue. When it doubles that cap, a \textbf{Patron Intrusion} occurs.

\section{Clearing Obligation – General Ways}
\begin{itemize}
    \item \textbf{Act of Service (Downtime):} Perform a task aligning with the Patron’s domain. Reduce Obligation by 1–2 (GM discretion).
    \item \textbf{Ritual Purification (Scene):} Spend a scene in prayer, meditation, or offering. Test \texttt{Lore + Spirit} (DV 3). Success reduces Obligation by 1. Failure may add a complication.
    \item \textbf{Patron’s Largess (GM Fiat):} When you dramatically fulfull the Patron’s unstated goal, the GM may clear all Obligation as a reward.
\end{itemize}

\section*{Patron Intrusions – Narrative Consequences}
When Obligation exceeds \texttt{2 × (Spirit+Presence)}, the GM triggers an intrusion. Below are examples. The GM should tailor them to the current scene.

% Patron Chapter – Template and Example
% Use within \chapter{Patrons of the Gray Wanderer} or similar

% ----------------------------------------------------------------------
% Template for a Patron Entry (to be copied for each patron)
% ----------------------------------------------------------------------
%
% \section{Patron Name – Epithet}
% \subsection*{Biography and Alignment}
% Flavor text / anecdote...
% 
% \textbf{What They Represent:} ...
% 
% \textbf{Cultural Variants:} ...
% 
% \subsection*{Playing a Follower of [Patron]}
% \textbf{Strategy:} ...
% 
% \textbf{Suggested Acts of Obligation:} ...
% 
% \subsection*{Rites – Strategic Use}
% \textbf{Low Rites:} ...
% \textbf{Standard Rites:} ...
% 
% \subsection*{Advanced Rites (High Tier)}
% \begin{patronbox}{[Rite Name]}
% \textbf{Cost:} ... \textbf{Effect:} ...
% \end{patronbox}
% 
% ----------------------------------------------------------------------

\section{Aveh – The Rider Behind the Storm}
\label{patron:aveh}

\subsection*{Biography and Alignment}

\textit{“I am the track the storm devours. I am the freedom you fear and the silence you crave. To ride with me is to be unbound – and to be unbound is to vanish.”}

Aveh is the faceless rider at the horizon’s edge, neither man nor woman, both and neither. The Ykrul call them \textit{Aveh}, a liminal spirit beyond hearth and vow, storm‑shadow that offers liberty at the cost of belonging. Wanderers, exiles, oath‑breakers, and rebels whisper Aveh’s name – envied for their freedom, condemned for their abandonment of duty.

\textbf{What They Represent:} Freedom, erasure, storms, exile. Aveh is the patron of those who have left behind everything – names, homes, even memories. They do not build; they unbind.

\textbf{Cultural Variants:}
\begin{itemize}
    \item \textbf{Ykrul:} A sky‑spirit that rides with winter gales; shamans invoke Aveh to hide a dying person’s tracks from the dead.
    \item \textbf{Linn:} A sea‑fog rider; sailors whisper Aveh’s name when they must slip through a blockade unseen.
    \item \textbf{Acasia:} The “dust‑walker” who appears at crossroads; followers leave a coin on the lintel when they flee a blood feud.
\end{itemize}

\subsection*{Playing a Follower of Aveh}

\textbf{Strategy:} Aveh’s rites excel at escape, stealth, and erasure. You are not a front‑line fighter – you are the one who slips away, who makes evidence vanish, who leaves no trace. In social scenes, you are forgettable by design.

\textbf{Suggested Acts of Obligation (to clear debt):}
\begin{itemize}
    \item Cut a binding contract (yours or another’s) and scatter the pieces to the wind.
    \item Help someone flee oppression – a slave, a debtor, a threatened child.
    \item Spend a night sleeping under open sky with no shelter, marking no camp.
    \item Sabotage a waypoint (bridge, toll gate, beacon) so that travellers must find a new path.
\end{itemize}

\subsection*{Rites – Strategic Use}

\textbf{Low Rites}
\begin{itemize}
    \item \textit{Storm‑Step} (4 XP): Slip free of notice or restraint for a scene. Use it to escape a grapple, lose a tail, or walk past a guard who just saw you.
    \item \textit{Exile’s Banner} (5 XP): Mark a target as “outside” – socially erased. Use on a rival before a negotiation to make their words carry no weight.
\end{itemize}

\textbf{Standard Rites}
\begin{itemize}
    \item \textit{Erase the Road} (8 XP): Trails vanish, records blur, pursuers lose the way. Perfect for after a heist or when hunted.
    \item \textit{The Rider’s Mark} (9 XP): Bestow Aveh’s mark on an ally – they resist capture but become forgettable. Use on a scout or messenger.
\end{itemize}

\subsection*{Advanced Rites (High Tier, 12+ XP)}

\begin{patronbox}{Storm’s Refuge}
\textbf{Cost:} 12 XP, Obligation 6 segments \\
\textbf{Effect:} Call on the Rider to shroud your company. For the rest of the scene, your group is unseen or forgotten by pursuers. At the end, all present must mark 1 Corruption or forget one bond or obligation. \\
\textbf{Strategic note:} Use when you must cross a guarded border or escape after a failed heist. The corruption cost is high – but sometimes silence is worth your soul.
\end{patronbox}

\begin{patronbox}{The Horizon Devours [BANE]}
\textbf{Cost:} 14 XP, Obligation 8 segments \\
\textbf{Effect:} The Rider swallows boundaries. Walls crumble, oaths dissolve, barriers fail. All thresholds in a zone collapse, but something of you is erased in turn – your reflection, your name, or your very presence in memory. \\
\textbf{Strategic note:} End a siege, break a prison, or nullify a magical ward. Use only when you are willing to become a ghost in your own story.
\end{patronbox}

\subsection*{GM Guidance – Aveh’s Intrusions}

When Aveh’s Obligation exceeds double capacity, choose or roll:
\begin{enumerate}
    \item A storm erases your tracks – and also the memory of your face from the last town you visited.
    \item Your shadow detaches and wanders off for an hour; when it returns, it brings a creepy omen.
    \item A forgotten debt suddenly resurfaces – someone you wronged now knows exactly where you are.
\end{enumerate}

\clearpage
% ----------------------------------------------------------------------
% Grimmir – The Wild Speaker (Druid Analogue)
% ----------------------------------------------------------------------
\section{Grimmir – The Wild Speaker}
\label{patron:grimmir}

\subsection*{Biography and Alignment}

\textit{“In the space between planted row and forest edge, where the first grain met the first acorn, Grimmir waits. He speaks in the rustle of leaves, the migration of birds, and the patient growth of ancient oaks. Listen, and learn that true power comes not from dominion, but from harmony.”}

Grimmir walks the threshold between cultivated lands and untamed wilderness. Neither wholly beast nor entirely human, he speaks in root and stone, seasons and sap. Those who walk his path gain insight into nature’s patterns, the tongues of leaf and fang, and mastery over the primal forces that shape the wild places – yet must choose which voice to trust: the cultivated mind or the feral heart.

\textbf{Vignette – A Huntress’s Confession:}
\epigraph{“I tracked a poacher for three days through the Valewood. On the second night, I heard him praying – not to any god, but to the trees themselves, begging forgiveness for each arrow he’d loosed. Grimmir didn’t answer. But the roots shifted, and by morning the man’s trail was gone. I reported failure. My captain called me a coward. The forest called me its own.”}{\textit{— Rina, Vilikari hedge-warden}}

\textbf{What They Represent:} Primal wisdom, seasonal cycles, the treaty between civilisation and wilderness. Grimmir is not a god of nature – he is nature’s patient accountant, balancing every cut tree with a new seedling, every hunt with a necessary death.

\textbf{Cultural Variants:}
\begin{itemize}
    \item \textbf{Aelaerem (Halflings):} Called \textit{the Thorn-Father}, keeper of hedgerow magic. His shrines are cup‑mark stones where butter is left for the Neighbors.
    \item \textbf{Ykrul:} A sky‑wolf who herds the seasons; shamans invoke Grimmir to predict the spring thaw or the first frost.
    \item \textbf{Vilikari:} A two‑faced figure – one side a sower, one side a reaper; he is invoked at both planting and harvest festivals.
    \item \textbf{The Inquisitor Prime’s View:} Witch‑hunters call Grimmir \textit{the Root Serpent}, a deceiver who tempts settlers to abandon their fields and return to savagery. His followers are hunted as heretics in Ecktoria.
\end{itemize}

\subsection*{Playing a Follower of Grimmir}

\textbf{Strategy:} Grimmir’s rites reward patience, observation, and environmental manipulation. You are a controller – slowing enemies with thorns, creating advantageous terrain, speaking with beasts to gather intelligence. In combat, you rarely strike directly; you make the battlefield itself into a weapon. Outside combat, you are the party’s guide, forager, and weather‑reader.

\textbf{Suggested Acts of Obligation (to clear debt):}
\begin{itemize}
    \item Plant a tree or restore a damaged grove (requires a scene of work; may be done during downtime).
    \item Kill a predator that has grown unnaturally bold and is upsetting the local balance.
    \item Teach a villager to read a seasonal sign (bird migration, flower bloom) – pass on knowledge rather than hoarding it.
    \item For one week, eat no meat that was not taken by your own hand or grown in your own plot.
\end{itemize}

\subsection*{Rites – Strategic Use}

\textbf{Low Rites}
\begin{itemize}
    \item \textit{The Speaking Seed} (4 XP): Commune with local plant life. Use to find water, identify poisonous flora, or ask a grove who last passed through.
    \item \textit{The Season’s Turn} (5 XP): Align with seasonal power. Gain +1 die to Survival in that season or resistance to its hazards. In spring, resist pollen‑induced sneezing (narrative, but your GM will appreciate it).
\end{itemize}

\textbf{Standard Rites}
\begin{itemize}
    \item \textit{Verdant Tongue} (8 XP): Speak with all plant and animal life in Near range. Perfect for interrogating a witness deer or negotiating with a territorial badger.
    \item \textit{The Thornveil} (7 XP): Raise a living barrier. Mark a zone as [WARD] against those who would harm the natural order. Use to protect a campsite or to funnel enemies into a kill box.
\end{itemize}

\subsection*{Advanced Rites (High Tier, 12+ XP)}

\begin{patronbox}{The World’s Wound}
\textbf{Cost:} 13 XP, Obligation 7 segments \\
\textbf{Effect:} Heal a blighted place or quicken a damaged ecosystem. Clear taint, restore fertility, awaken dormant forces. Gain +2 dice to related Survival/Nature actions in the region. Begin an \textit{Ecosystem Restoration} [8] clock. \\
\textbf{Strategic note:} Use after a magical disaster (e.g., a Khemesh corruption) or to earn the gratitude of a forest community. The clock tracks the recovery – if you abandon it halfway, the wound festers worse than before.
\end{patronbox}

\begin{patronbox}{The Cycle’s Crown}
\textbf{Cost:} 14 XP, Obligation 8 segments \\
\textbf{Effect:} Become an avatar of the seasonal cycle for one scene. Choose two: +2 dice to all actions keyed to the current season; command plant life in Near automatically; remove one Condition or downgrade Harm by one level; allies gain +1 Armor from natural cover. Begin a \textit{Seasonal Avatar} [6] clock. \\
\textbf{Strategic note:} This is your “boss fight” mode. Use it when defending a sacred grove or when outnumbered in deep wilderness. When the clock fills, the season shifts dramatically (GM’s choice) – maybe an early frost or a sudden thaw.
\end{patronbox}

\subsection*{GM Guidance – Grimmir’s Intrusions}

When Grimmir’s Obligation exceeds double capacity, choose or roll:
\begin{enumerate}
    \item A tree you passed yesterday is now blocking your path – with a face carved into it that looks like yours.
    \item All food you touch spoils within an hour unless you offer the first bite to the earth.
    \item A pack of wolves shadows the party for a full day. They do not attack – they simply watch, judging.
\end{enumerate}

\noindent \textbf{Witch‑Hunter Variant (Inquisitor Prime’s Perspective):}
If a PC follows Grimmir in a region hostile to wild magic, the GM may replace the above intrusions with \textit{Inquisitorial Pressure}: local officials plant cursed iron nails at crossroads, farmers refuse to sell to the party, and a bounty poster appears with a leaf‑shaped brand.

\clearpage
\endinput

% ----------------------------------------------------------------------
% Ikasha – She Who Sleeps (Latent Potential & Shadow)
% ----------------------------------------------------------------------
\section{Ikasha – She Who Sleeps}
\label{patron:ikasha}

\subsection*{Biography and Alignment}

\textit{“Blow out the candle. If the room listens back, ask softly. At the next crossroads, the raven waits – and the shadow remembers your passing.”}

Ikasha is the hush between footfalls, the patience of dark water, the black‑feathered watcher at every threshold. In stillness she gathers what might be, in crossroads she whispers of what may yet come. Ravens circle her, bearing secrets between worlds. Her followers learn to move unseen and speak unremembered, becoming shadows that slip between what is and what could be.

\textbf{Vignette – A Runekeeper’s Regret:}
\epigraph{“I bound a Rite of Ikasha into my Codex on a dare. For a year, it was useless – a blank page that rustled when no wind blew. Then I was cornered in an alley by three Inquisitorial agents. I whispered, ‘Sleep,’ and the page tore itself free. The ravens came. When I woke, the alley was empty. So was the page. So was my memory of the last hour.”}{\textit{— Vessa of Silkstrand, exiled Runekeeper}}

\textbf{What They Represent:} Latent potential, shadow, thresholds, forgotten things. Ikasha does not act; she \textit{waits}. Her power is the knife not yet drawn, the word not yet spoken, the path not yet taken. Followers learn that patience is sharper than any blade.

\textbf{Cultural Variants:}
\begin{itemize}
    \item \textbf{Mistlands:} Called \textit{the Bell‑Shrouded}; bell‑wardens leave a single rope unknotted for her, so she may ring silence when needed.
    \item \textbf{Zakov:} A smuggler’s patron – \textit{the Tide Watcher}. She rises with the night tide, and her ravens carry messages between pirate captains.
    \item \textbf{Aeler:} \textit{the Deep Listener}; miners whisper to her before descending into new shafts, hoping she will keep the stone from speaking their names aloud.
\end{itemize}

\subsection*{Playing a Follower of Ikasha}

\textbf{Strategy:} Ikasha’s rites reward stealth, deception, and strategic patience. You are not a brawler – you are the one who waits for the guard’s shift change, who slips through the door left ajar, who knows that a secret kept is worth ten swords drawn. In social scenes, you excel at staying unnoticed or, when needed, becoming unforgettable in the worst way.

\textbf{Suggested Acts of Obligation (to clear debt):}
\begin{itemize}
    \item Spend a night in absolute darkness, speaking no word and lighting no flame.
    \item Carry a secret for someone – not a small one, but a truth that would ruin them – and tell no one, not even your party, for a full season.
    \item At a crossroads, leave a black feather and a single coin. Do not look back.
    \item Sabotage a lantern or torch that would have revealed someone hiding.
\end{itemize}

\subsection*{Rites – Strategic Use}

\textbf{Low Rites}
\begin{itemize}
    \item \textit{Touch the Umbral Veil} (4 XP): Start Controlled on one Stealth roll or gain +1 Effect to hide/move quietly. Use when you absolutely must not be seen – the first roll of a heist.
    \item \textit{Rite of the Crossroads Raven} (5 XP): Summon an omen‑raven. Grant +1 die to a navigation, pursuit, or diversion action, or force an enemy to hesitate. The raven can also carry a whispered message up to Far range.
\end{itemize}

\textbf{Standard Rites}
\begin{itemize}
    \item \textit{Draw from the Umbral Reservoir} (8 XP): Grant yourself or an ally +2 dice to stealth, deception, or resolve, or clear 1 Fatigue. Use when someone needs to shake off the mental strain of a failed infiltration.
    \item \textit{Secret Keeper’s Burden} (9 XP): Compel a truthful answer to one direct question. Deep secrets allow a Resolve test to resist. Perfect for interrogating a captured spy or extracting a confession.
\end{itemize}

\subsection*{Advanced Rites (High Tier, 12+ XP)}

\begin{patronbox}{Become the Shadow at the Crossroads}
\textbf{Cost:} 12 XP, Obligation 7 segments \\
\textbf{Effect:} You become intangible to mundane harm, pass through thresholds and small gaps, gain +2 dice to Stealth, and auto‑succeed one escape attempt. However, you cannot manipulate normal objects. \\
\textbf{Strategic note:} Use to bypass a locked door, escape a collapsing building, or spy on a conversation through a keyhole. The ritual ends if you try to interact with the physical world – so plan your exit before you phase in.
\end{patronbox}

\begin{patronbox}{The Raven’s Debt [BANE]}
\textbf{Cost:} 14 XP, Obligation 8 segments \\
\textbf{Effect:} Name a target you have seen within the past day. Three ravens will stalk them relentlessly for one week. The target is never alone, never unobserved, and suffers –2 dice to all stealth and deception rolls. Once per day, the GM may have the ravens reveal one secret the target tried to hide. \\
\textbf{Strategic note:} This is not assassination – it is psychological warfare. Use on a rival who thinks they have outmanoeuvred you. The ravens cannot be killed by mundane means; they return an hour later. The only way to end the curse is to apologise publicly to the follower of Ikasha who cast it.
\end{patronbox}

\subsection*{GM Guidance – Ikasha’s Intrusions}

When Ikasha’s Obligation exceeds double capacity, choose or roll:
\begin{enumerate}
    \item A raven lands on your shoulder and whispers a secret you never wanted to know – about an ally, a foe, or yourself.
    \item Your shadow stretches into the shape of a woman with a crown of black feathers. For the next hour, anyone who sees your shadow must make a Resolve test (DV 3) or become convinced they are being followed.
    \item A door you passed through earlier is now locked from the other side – and the key is in your pocket, even though you never picked it up.
\end{enumerate}

\noindent \textbf{Cantor Variant – The Silence Between Notes:}
A Cantor of Ikasha does not sing loudly. She hums at the edge of hearing, her melodies felt more than heard. Her \textit{Resonant Rite} is a pause – a rest in the music where the audience leans forward, and in that leaning, forgets what they were doing. Use Performance Maneuver (Position shift) to represent this; the sound is so subtle that enemies don’t realise they’ve lost their nerve.

% ----------------------------------------------------------------------
% Khemesh – The Abyssal Maw (Eldritch / Deep Horror)
% ----------------------------------------------------------------------
\section{Khemesh – The Abyssal Maw}
\label{patron:khemesh}

\subsection*{Biography and Alignment}

\textit{“In the trench without light, the Maw waits. Even silence drowns.”}

Khemesh is not merely a lord of the depths – he is the hunger beneath them, a pressure older than seas. Those who bargain with him are marked by the abyss, seen in the way shadows cling, in the whispers heard when no voice speaks, in the certainty that all things will sink. He is son to Raéyn, the Mistress of the Sea, and embodies her greatest tragedy: the crushing inevitability of the ocean’s dark heart.

\textbf{Vignette – A Cantor’s Final Hymn:}
\epigraph{“We found the Cantor on the shore, still clutching her conch shell. Her eyes were open, but they saw nothing. Her lips moved, but no sound came. The sailors said she had sung a Rite of Khemesh to calm a storm – and the storm obeyed. But when it passed, it left something in her. Something that now waits in the deep, wearing her voice.”}{\textit{— Log of the Sea‑Witch, last entry}}

\textbf{What They Represent:} Depth, pressure, inevitability, eldritch terror. Khemesh is the final accounting – not evil, but undeniable. He does not punish; he \textit{applies weight}. His followers learn that resistance is vanity, and that the only freedom is to sink on your own terms.

\textbf{Cultural Variants:}
\begin{itemize}
    \item \textbf{Zakov:} \textit{the Drowned Prince}. Pirates toss a silver coin overboard before a raid – if it sinks, Khemesh approves. If it floats, they turn back.
    \item \textbf{Linn:} \textit{the Reef‑Father}. Fishermen leave the first catch of the day to him, lest the tide turn and drag their nets into the abyss.
    \item \textbf{Mistlands:} \textit{the Bell‑Breaker}. When a bell line fails, the wardens whisper that Khemesh has breathed on the clappers. Repairs are made with iron from a sunken ship.
\end{itemize}

\subsection*{Playing a Follower of Khemesh}

\textbf{Strategy:} Khemesh’s rites excel at control, crushing force, and psychological terror. You are a hammer – slow, inevitable, and devastating. In combat, you pin enemies, create zones of crushing pressure, and erode enemy morale. Outside combat, you are an investigator of the deep, reading secrets from water and stone.

\textbf{Suggested Acts of Obligation (to clear debt):}
\begin{itemize}
    \item Throw a valuable object into the deepest water you can reach (a lake, a well, the sea). It must have personal meaning.
    \item Spend a night submerged up to your neck in cold water (river, ocean, bathhouse). Mark 1 Fatigue – this is a meditation, not a punishment.
    \item Reveal a secret you have kept for more than a year. Speak it aloud where water can hear.
    \item Repair a broken bell or a fractured levee – Khemesh respects those who mend boundaries he has tested.
\end{itemize}

\subsection*{Rites – Strategic Use}

\textbf{Low Rites}
\begin{itemize}
    \item \textit{Whisper of the Trench} (4 XP): Target suffers –1 die on their next action from disorienting echoes of the deep. Use to soften a boss before an ally’s big strike.
    \item \textit{Rite of Crushing Silence} (5 XP): Establish an oppressive silence in a zone. Enemies suffer –1 die to coordination and morale actions. Use to disrupt enemy commanders or to sneak past guards who rely on call‑and‑response.
\end{itemize}

\textbf{Standard Rites}
\begin{itemize}
    \item \textit{Pressure of the Maw} (7 XP): Pin a target with invisible force – treat as [ENTANGLE] with Great Effect if underwater or confined. Use to immobilise a fleeing foe or to hold a door shut.
    \item \textit{Rite of the Abyssal Vision} (9 XP): See the world as Khemesh does – fractured, alien, crushing. Gain +2 dice to Notice and Arcana, and ask one “true nature” question about a foe or structure. The cost is Exposure +1 when the scene ends (your eyes take time to readjust).
\end{itemize}

\subsection*{Advanced Rites (High Tier, 12+ XP)}

\begin{patronbox}{The Maw Opens}
\textbf{Cost:} 12 XP, Obligation 8 segments \\
\textbf{Effect:} Reality in the zone folds inward like the crushing deep. Enemies act at Desperate Position by default. Each beat, the Keeper may force 1 SB (Spades/Clubs favoured) as the pressure mounts. Structures, vessels, or wards fracture as if under immense weight. \\
\textbf{Strategic note:} Use this to end a siege or to break a fortified position. The zone affects everyone, including allies – so cast it, then pull back and let the enemy drown in their own shattered morale.
\end{patronbox}

\begin{patronbox}{The Inevitable Descent [CURSE]}
\textbf{Cost:} 14 XP, Obligation 8 segments \\
\textbf{Effect:} Name a target you have wounded in the past minute. For the rest of the scene, that target cannot escape the zone by any means – no teleportation, no flight, no hidden passages. They also suffer –2 dice to all movement‑related rolls. At the end of the scene, they must make a Spirit+Resolve test (DV 4) or take Harm 2 (crushing). \\
\textbf{Strategic note:} This is your “no escape” button. Use it on a slippery rogue, a flying enemy, or anyone who thinks they can outrun justice. The curse ends when the target yields or is defeated – but yield must be accepted.
\end{patronbox}

\subsection*{GM Guidance – Khemesh’s Intrusions}

When Khemesh’s Obligation exceeds double capacity, choose or roll:
\begin{enumerate}
    \item Water pressure quadruples in the room – all wooden objects groan, every breath feels heavy. Glassware cracks.
    \item A single nightmare tentacle manifests from a nearby body of water (or from the floor if none exists) and steals a shiny object from the party – then retreats.
    \item The party begins to \textit{sink} – not physically, but emotionally. Each character must mark 1 Fatigue or confess a secret fear aloud.
\end{enumerate}

\noindent \textbf{Runekeeper Note:} Khemesh’s Codex is a **leather‑bound log of shipwrecks**. Each Rite inscribed requires a drop of salt water on the page. If the Codex dries out completely, it becomes unreadable until soaked in sea water for a day.

% ----------------------------------------------------------------------
% Carrion‑King – Decay, Renewal & Transformation (Dark Druid / Death Druid)
% ----------------------------------------------------------------------
\section{Carrion‑King – The Master of Endings}
\label{patron:carrionking}

\subsection*{Biography and Alignment}

\textit{“What crumbles feeds what grows. What dies becomes the soil of tomorrow’s triumph. Do not mourn the fallen – harvest them.”}

The Carrion‑King is not a lord of rot for its own sake. He is the alchemist of endings, the patient gardener who sees every corpse as compost, every failed empire as fertiliser for the next. His followers are pragmatists who understand that destruction is simply the first step of creation. They are not cruel – they are *efficient*.

\textbf{Vignette – An Invoker’s Bargain:}
\epigraph{“I knelt in the plague pit, my Symbol of the Carrion‑King pressed to my chest. The dying soldier coughed blood on my boots. ‘Save me,’ he whispered. I could not. Instead, I asked the King to let the man’s death feed the village’s winter wheat. The King agreed. The wheat grew thick and black, but it fed the children. The soldier’s name is carved on the granary door.”}{\textit{— Isa of Acasia, Hedge‑Invoker}}

\textbf{What They Represent:} Decay, renewal, transformation. The Carrion‑King is the principle that nothing is wasted – only repurposed. He is worshipped by farmers, gravediggers, and those who have lost everything and must rebuild.

\textbf{Cultural Variants:}
\begin{itemize}
    \item \textbf{Vhasia:} \textit{the Dust‑Accountant}. In the war‑torn marches, soldiers leave a button from a fallen comrade on a cairn – the King will turn that button into a bullet for the next battle.
    \item \textbf{Aeler:} \textit{the Under‑Gardener}. Dwarven crypts have a patch of bioluminescent fungus fed by the dead. Tending it is an act of devotion.
    \item \textbf{The Inquisitor Prime’s View:} Witch‑hunters see the Carrion‑King as a necromantic menace, confusing ‘natural decay’ with ‘unholy resurrection’. His followers are burned with the bodies they tried to compost.
\end{itemize}

\subsection*{Playing a Follower of the Carrion‑King}

\textbf{Strategy:} You are a harvester of potential. In combat, you weaken enemies with rot and decay, then use their defeat to empower your allies. Out of combat, you turn obstacles into opportunities – a collapsed bridge becomes timber, a dead horse becomes leather and glue. You are the party’s logistician and recycler.

\textbf{Suggested Acts of Obligation (to clear debt):}
\begin{itemize}
    \item Compost a fallen enemy’s body (or a slain animal) and use the soil to plant something edible.
    \item Salvage a ruined structure – a collapsed mill, a burned barn – and give the reclaimed materials to a needy family.
    \item Teach someone that loss is not failure. Spend a scene counselling a grieving villager.
    \item Burn a written lie (a false contract, a slanderous letter) and scatter the ashes on fallow ground.
\end{itemize}

\subsection*{Rites – Strategic Use}

\textbf{Low Rites}
\begin{itemize}
    \item \textit{Rite of Consuming Rot} (5 XP): Accelerate natural decay to weaken organic materials (ropes, leather, wood). Gain +2 Effect to Break/Sabotage against such targets. Use to snap a bridge rope or rot a leather armour strap before a fight.
    \item \textit{Rite of the Harvested End} (4 XP): Extract value from a recent ending – a defeated foe, a failed plan, a broken item. Gain +1 die to your next action, re‑roll a 1, or gain 1 SB to spend immediately. Use to turn a failure into fuel.
\end{itemize}

\textbf{Standard Rites}
\begin{itemize}
    \item \textit{Rite of the Fertile Death} (8 XP): Transform death into growth – create advantageous terrain (cover, concealment) or grant allies +1 die to healing/recovery rolls. Use after a battle to turn a bloodied field into a temporary refuge.
    \item \textit{Rite of the Transformed Spirit} (7 XP): Channel the essence of a recently deceased being. Gain one skill die reflecting their expertise for a scene, or ask one question about their knowledge. Perfect for extracting information from a dead enemy’s spirit.
\end{itemize}

\subsection*{Advanced Rites (High Tier, 12+ XP)}

\begin{patronbox}{The Great Consumption}
\textbf{Cost:} 13 XP, Obligation 7 segments \\
\textbf{Effect:} Transform a large area through decay and renewal. Choose two: create difficult terrain that favours you and your allies, summon a Cap 3 swarm of scavengers as temporary allies (ravens, rats, beetles), or generate valuable reagents worth 2 XP (dried herbs, bone meal, rare fungi). \\
\textbf{Strategic note:} Use this after you have secured a battlefield but before you leave. The scavengers can act as a rear guard, the terrain slows pursuers, and the reagents pay for your next downtime. A rite of practical, not flashy, power.
\end{patronbox}

\begin{patronbox}{The Eternal Cycle [TRANSFORM]}
\textbf{Cost:} 14 XP, Obligation 8 segments \\
\textbf{Effect:} Complete a transformation cycle. Destroy one major asset, enemy, or obstacle, and create something new of equal or greater value. GM and player collaborate to define the transformation – e.g., a cursed grove becomes a fertile orchard; a noble’s stolen treasure map becomes a key to a different vault. \\
\textbf{Strategic note:} This is a campaign‑level rite. Use it when a situation has reached a dead end and you need to *change the terms*. The Carrion‑King does not allow pure destruction – you must build as you break.
\end{patronbox}

\subsection*{GM Guidance – Carrion‑King’s Intrusions}

When the Carrion‑King’s Obligation exceeds double capacity, choose or roll:
\begin{enumerate}
    \item A dead creature you passed earlier is now following you – not as an enemy, but as a witness. It will not attack, but it will not leave.
    \item All food in the party’s packs spoils instantly, except for a single loaf of black bread that tastes of ash but will not rot. It sustains you for a week, but you must eat it alone.
    \item A villager you saved from death now bears a small rot‑mark on their skin. It does not hurt, but it spreads one inch each day. The only cure is to complete an act of service for the King.
\end{enumerate}

\noindent \textbf{Cantor Variant – The Dirge of Compost:}
Cantors of the Carrion‑King sing low, humming melodies that sound like buzzing flies. Their Performance Maneuver does not inspire courage – it inflicts a –1 Position penalty on enemies who hear it (instead of buffing allies). Choose at casting.

% ----------------------------------------------------------------------
% Morag the Hag – Weaver of Hidden Costs (Faerie Bargains & Hidden Costs)
% ----------------------------------------------------------------------
\section{Morag the Hag – Weaver of Hidden Costs}
\label{patron:morag}

\subsection*{Biography and Alignment}

\textit{“She gives you the gold, but not the weight it carries. She grants your heart’s desire, but not the hunger it awakens. She answers your plea, but not the price that pleases her. In Morag’s court, every ‘yes’ is a ‘yes, but…’”}

In the hour between twilight and dawn, where borders thin and mortal cunning meets older cleverness, Morag sets her webs of bargain and consequence. She is the grandmother who offers exactly what you need, the forest crone with a perfect solution, the crossroads dweller who makes dreams come true – for a price that only reveals itself once the road bends back. Her wisdom is simple and sharp: every gift carries obligation, every kindness creates debt, every wish must be paid in coin not named aloud.

\textbf{Vignette – A Runekeeper’s Regret:}
\epigraph{“I borrowed a horse from a stranger at a crossroads. She said, ‘Return it by dawn, and we’re square.’ I did. The horse was fine. But the stable boy had vanished. I later learned the horse was him – transformed for a night’s labour. I had returned a horse, not a child. Morag kept the difference.”}{\textit{— Runekeeper Torvin, now a goat herder}}

\textbf{What They Represent:} Faerie bargains, hidden costs, old‑road wisdom. Morag is not evil – she is *precise*. She respects the letter of a deal, but the spirit is hers to interpret. Followers learn to read between the lines, to count the threads, and to never, ever accept a gift without asking, “What does this cost the giver?”

\textbf{Cultural Variants:}
\begin{itemize}
    \item \textbf{Aelaerem (Halflings):} \textit{the Hearth‑Grandmother}. She is invoked before large family negotiations, and her symbol is a knotted red thread. Children leave a crust of bread on the windowsill to keep her from “improving” their toys.
    \item \textbf{Acasia:} \textit{the Bridge‑Wife}. At toll bridges, a small side‑shrine offers her a coin in exchange for safe passage – but the coin must be one you found, not one you earned.
    \item \textbf{Valewood:} \textit{the Bark‑Scrivener}. Fae courts keep a copy of every bargain witnessed by Morag. Altering a contract after her seal is impossible – the paper will burn itself.
\end{itemize}

\subsection*{Playing a Follower of Morag}

\textbf{Strategy:} You are the party’s negotiator, trap‑layer, and loophole finder. Morag’s rites excel at binding agreements, creating escape clauses, and trading disadvantage for advantage. In combat, you are a trickster – you do not fight fair; you make the enemy trip over their own feet, forget their orders, or suffer their own reflected hubris.

\textbf{Suggested Acts of Obligation (to clear debt):}
\begin{itemize}
    \item Complete a bargain where you receive less than you give – intentionally take the worse end of a deal.
    \item Return a “found” object to its rightful owner, even if you have no legal obligation to do so.
    \item Offer hospitality to a stranger for a full night (food, fire, shelter). If they steal from you, forgive the debt.
    \item Burn a contract you have signed, then scatter the ashes at a crossroads at midnight.
\end{itemize}

\subsection*{Rites – Strategic Use}

\textbf{Low Rites}
\begin{itemize}
    \item \textit{Rite of the Unfinished Promise} (4 XP): Frame a bargain with an escape clause only you understand. Gain +1 die when the other party acts to fulfil their obligation. Use when making a deal with an untrustworthy NPC – the die bonus applies when they try to betray you.
    \item \textit{Rite of the Borrowed Luck} (5 XP): Share in another’s fresh success. Gain +2 dice on one roll that benefits from their good fortune. Use immediately after an ally succeeds on a crucial roll – you are riding their wave.
\end{itemize}

\textbf{Standard Rites}
\begin{itemize}
    \item \textit{Rite of the Speaking Knot} (8 XP): Bind a promise with fae knots that tighten when broken. Both parties gain +1 die while keeping the bargain; breaking it escalates consequences (Harm 1, then Harm 2, then Curse). Perfect for securing a fragile alliance.
    \item \textit{Rite of the Threefold Exchange} (7 XP): Broker a trade where each party yields hidden value. All participants gain +1 die when the exchange benefits them; Morag claims the difference in true worth as her tithe. Use when you want to mediate a dispute – everyone feels they gained, but you know Morag took her cut.
\end{itemize}

\subsection*{Advanced Rites (High Tier, 12+ XP)}

\begin{patronbox}{The Cauldron’s Secret}
\textbf{Cost:} 13 XP, Obligation 7 segments \\
\textbf{Effect:} Brew a solution to any problem – a draught that heals a plague, a poultice that wakes a sleeper, a poison that only affects the guilty. The cauldron asks one more ingredient: one the brewer does not realise they are providing until too late (e.g., a cherished memory, a drop of blood from a loved one, a year of your life). \\
\textbf{Strategic note:} Use only when the alternative is certain death or failure. The ingredient is always something you cannot afford to lose – but you will only discover what *after* the brew is complete. Morag loves irony.
\end{patronbox}

\begin{patronbox}{The Crossroads Court [BIND]}
\textbf{Cost:} 14 XP, Obligation 8 segments \\
\textbf{Effect:} Convene Morag’s court at a true crossroads at twilight. All present must bargain. Each pact grants a marked boon (e.g., +1 die to all rolls in a coming battle) and hides a cost that will ripen at the worst time (e.g., your hair turns white, an old rival reappears). The court lasts until dawn; leaving early breaks all bargains and summons a minor fae enforcer. \\
\textbf{Strategic note:} This is a session‑long set piece. Use it to force a reluctant ally to commit, to trick an enemy into a losing deal, or to resolve a multi‑party standoff. The cost is always narrative, never purely mechanical – but it will hurt.
\end{patronbox}

\subsection*{GM Guidance – Morag’s Intrusions}

When Morag’s Obligation exceeds double capacity, choose or roll:
\begin{enumerate}
    \item A coin you spent yesterday is back in your pocket today – but now it has a single strand of red thread glued to it.
    \item An NPC who owes you a favour now claims you owe *them* – they produce a folded note in your handwriting that you do not remember writing.
    \item The party’s food turns into a feast (plentiful, delicious) – but every dish is slightly wrong: the bread is hollow, the wine is coloured water, the meat is a mushroom. You are not poisoned; you are merely *unsatisfied*.
\end{enumerate}

\noindent \textbf{Cantor Variant – The Hum of the Unsaid:}
Cantors of Morag sing in round – overlapping melodies that create a third, hidden tune. Their Performance Maneuver does not shift ally Position; instead, it delays an enemy’s action by one beat (they must re‑roll their initiative or act last in the exchange). The audience never remembers the melody, only the feeling of being *out‑manoeuvred*.

\clearpage

% ----------------------------------------------------------------------
% Nidhoggr – The World-Worm (Dreaming Antiquity & Stone Memory)
% ----------------------------------------------------------------------
\section{Nidhoggr – The World-Worm}
\label{patron:nidhoggr}

\subsection*{Biography and Alignment}

\textit{“Press your ear to the earth and wait. If it remembers you, it will answer.”}

Beneath stone and sleep coils Nidhoggr, who gnaws at the roots of time. He does not speak quickly; he dreams in centuries. To press your ear to the earth is to risk drowning in the silence of aeons. Yet for those who endure, he whispers truths long buried, memories fossilised in stone, and the slow inevitability of cycles unbroken. His followers walk in twilight between dream and ruin, bearing the weight of all that has been.

\textbf{Vignette – A Geomancer’s Discovery:}
\epigraph{“I was surveying a collapsed aqueduct near Khaz‑Vurim when my rod struck something soft – not earth, not stone. I pressed my ear to the crack and heard breathing. Slow. Rhythmic. Centuries long. I whispered, ‘Who sleeps?’ The earth answered: ‘Everyone you ever knew. And you will join them.’ I do not survey anymore.”}{\textit{— Durin, former Aeler engineer}}

\textbf{What They Represent:} Dreaming antiquity, fossilised memory, geological time, the patience of stone. Nidhoggr is not a god of the past – he *is* the past, compressed and still alive. His followers are historians, archivists, and explorers who understand that every ruin has a heartbeat.

\textbf{Cultural Variants:}
\begin{itemize}
    \item \textbf{Aeler (Dwarves):} \textit{the Stone‑Father}. Miners leave a pinch of their own hair in a new shaft’s first crack – Nidhoggr will remember they were there.
    \item \textbf{Valewood:} \textit{the Root‑Seer}. The Lethai‑ar commune with Nidhoggr through petrified trees, reading the rings for ancient prophecies.
    \item \textbf{Theona:} \textit{the Ninth Bell}. Some well‑keepers whisper that Nidhoggr is the source of the “no ninth” taboo – he is the unspoken count, the thing that should not be measured.
\end{itemize}

\subsection*{Playing a Follower of Nidhoggr}

\textbf{Strategy:} You are the party’s memory – of places, of lore, of dangers past. Nidhoggr’s rites excel at investigation, prophecy, and slow, inevitable power. In combat, you are a controller who reshapes terrain, reads enemy weaknesses, and imposes the weight of ages. Out of combat, you are the scholar who finds the forgotten door, the archivist who knows the old name.

\textbf{Suggested Acts of Obligation (to clear debt):}
\begin{itemize}
    \item Restore a broken or defaced monument – a standing stone, a statue, a milestone.
    \item Retrieve an object that has been buried for more than a century and return it to its rightful place (not to a museum – to the earth).
    \item Spend a full night sleeping on bare stone. Mark 1 Fatigue, but gain a prophetic dream (GM answers one yes/no question about the next session).
    \item Teach someone a forgotten history – not a grand epic, but a small truth: the name of a street’s original builder, the reason a well was dug.
\end{itemize}

\subsection*{Rites – Strategic Use}

\textbf{Low Rites}
\begin{itemize}
    \item \textit{Stone’s Whisper} (4 XP): Stone reveals one hidden fact about its history. Use to learn who carved a secret door, when a ruin collapsed, or what was buried beneath the floor.
    \item \textit{Dreaming Deep} (5 XP): Enter a brief trance to witness echoes of local history. Ask one specific question about past events here. Use when you need to solve a mystery – the stone will show you what happened, not who did it.
\end{itemize}

\textbf{Standard Rites}
\begin{itemize}
    \item \textit{Stone‑Sleeper’s Murmur} (8 XP): Read accumulated memories in stone. Ask up to three questions about significant events witnessed here. Perfect for crime scene investigation or exploring an ancient tomb.
    \item \textit{Awakened Chronicle} (9 XP): Area replays ghostly echoes of a chosen past event. All witnesses gain +2 dice to Investigation/Recall related checks for the scene. Use when you need to walk a group through a historical re‑enactment – or trap an enemy in a loop of a defeat they have already suffered.
\end{itemize}

\subsection*{Advanced Rites (High Tier, 12+ XP)}

\begin{patronbox}{The World‑Worm’s Dream}
\textbf{Cost:} 13 XP, Obligation 7 segments \\
\textbf{Effect:} Access deep geological memory. Ask three questions about ancient history or receive a prophetic dream that covers a full session’s events (the GM will give you three cryptic statements that will prove accurate). You must spend the night in physical contact with stone – a cave, a quarry, a large boulder. \\
\textbf{Strategic note:} Use before a major expedition into unknown territory. The answers are always true, but they are true from the perspective of *stone* – which does not measure time the way you do. “The enemy sleeps beneath your feet” might mean a trap floor, a buried monster, or a metaphor.
\end{patronbox}

\begin{patronbox}{Aeon’s Eclipse [BANE]}
\textbf{Cost:} 14 XP, Obligation 8 segments \\
\textbf{Effect:} Overlay past reality onto the present location. Ruins partially reform, ghostly figures manifest, and the terrain reverts to a previous geological state (e.g., a valley becomes a river, a hill becomes a coastal cliff). Allies gain +2 dice to Lore/Investigation; modern technology or recent structures suffer –2 dice to function. The effect lasts until the next lunar cycle or until a named condition is met. \\
\textbf{Strategic note:} This is a campaign‑level rite. Use it to reveal a buried entrance, to drive out squatters who cannot abide the ghostly patrols, or to claim an ancient treasure that exists only in the past. However, the past may not want to let you leave.
\end{patronbox}

\subsection*{GM Guidance – Nidhoggr’s Intrusions}

When Nidhoggr’s Obligation exceeds double capacity, choose or roll:
\begin{enumerate}
    \item The ground beneath your feet turns to chalk. Your footprints remain – but they are from a different season, preserved as if millions of years old.
    \item You hear a voice you recognise – a dead relative, a lost friend – speaking from a stone wall. They are not trapped; they are merely *remembered*.
    \item A significant object (a key, a weapon, a letter) is now fossilised. It is still usable, but weighs three times as much and carries the scent of ancient seabed.
\end{enumerate}

\noindent \textbf{Runekeeper Note:} Nidhoggr’s Codex is carved onto stone tablets, not bound in leather. Carrying it requires a wagon or a very strong back. To cast a Rite, you must physically touch the tablet – so you cannot hide it in a pack. This is intentional; Nidhoggr does not favour the secretive.

\clearpage


% ----------------------------------------------------------------------
% Oath of Flame & Light – The Dawnfire (Paladin / Crusader)
% ----------------------------------------------------------------------
\section{Oath of Flame \& Light – The Dawnfire}
\label{patron:oathflame}

\subsection*{Biography and Alignment}

\textit{“Swear in the light. Keep it, or the light will keep you.”}

The Oath of Flame & Light is no patron of half‑measures. Its fire names, binds, and burns – demanding that those who swear within its radiance stand openly, speak truly, and pay the cost of keeping their word. At dawn altars, the sworn kindle sparks of consecrated fire; in battle, they blaze as torches that hold back the night. To follow this Oath is to live in public truth, with no shadow to hide in and no retreat from the vow once spoken.

\textbf{Vignette – A Cantor’s Ashes:}
\epigraph{“She sang the Hymn of Unbroken Promise at the siege of Tor Garak – a song so bright that the archers on the walls wept and dropped their bows. The enemy captain offered her gold to stop. She refused. He offered her safe passage. She refused. Then he offered to spare the village if she would only say, ‘I yield.’ She sang louder. The next morning, we found her. Her mouth was still open, and the coals around her were still warm. The village stood.”}{\textit{— Varic, Vilikari banner‑bearer}}

\textbf{What They Represent:} Sacred vows, divine justice, radiant power, dawn’s hope. The Oath is not a god – it is a covenant, a promise witnessed by the sun itself. Followers are paladins, templars, and oath‑bound knights who believe that some things are worth dying for, and that the light remembers every sacrifice.

\textbf{Cultural Variants:}
\begin{itemize}
    \item \textbf{Ecktoria (Everflame):} \textit{the Unquenchable Fire}. The state religion; priests carry censers of perpetual coal. Their Rites focus on purification and judgment.
    \item \textbf{Vilikari (Adar – the Dawnfire):} A war‑god of vows and vengeance. Invoked before battle, and after, to witness that no oath was broken.
    \item \textbf{Aeler:} \textit{the Light of Forges}. Dwarven smiths kindle their hearths with a brand from the temple’s eternal flame. The Oath here is less about justice and more about *craft integrity* – a flawed sword violates the covenant.
    \item \textbf{The Inquisitor Prime’s View:} The Inquisitor Prime and the Oath are frequent allies against corruption, but not the same. The Oath demands *personal* purity; the Inquisitor demands *absolute* purity. Where they conflict, schisms burn bright.
\end{itemize}

\subsection*{Playing a Follower of the Oath}

\textbf{Strategy:} You are the party’s shield and moral compass. The Oath’s rites excel at protection, healing, revealing lies, and punishing oath‑breakers. In combat, you are a front‑liner who can smite enemies, ward allies, and hold the line. Out of combat, you are a judge, a truth‑speaker, and a light in dark places – but your rigidity can alienate those who live in grey.

\textbf{Suggested Acts of Obligation (to clear debt):}
\begin{itemize}
    \item Swear a new vow – aloud, in front of witnesses – to protect someone or accomplish a task. Keep it, obviously.
    \item Offer shelter to a stranger for a night, asking nothing in return. If they betray your hospitality, forgive them once.
    \item Cleanse a defiled holy site (a desecrated shrine, a corrupted grove). Spend a scene in prayer and manual labour.
    \item Publicly admit to a lie you have told, no matter how small. The light demands truth.
\end{itemize}

\subsection*{Rites – Strategic Use}

\textbf{Low Rites}
\begin{itemize}
    \item \textit{Kindle Vow} (4 XP): Declare a short vow for this scene (e.g., “protect the villagers,” “speak only truth”). Bearer gains +1 die to actions fulfilling it. Use before a critical skill test where you need a boost.
    \item \textit{Lay on Hands} (5 XP): Cleanse affliction, downgrade Harm by 1, or remove 1 Fatigue. Use to keep a wounded ally in the fight or to purge a poison before it spreads.
\end{itemize}

\textbf{Standard Rites}
\begin{itemize}
    \item \textit{Sunlit Parley} (8 XP): Establish terms in open light. Honest persuasion gains +1 die; deceit suffers –1 die. Use when negotiating with someone you suspect is lying – the light makes it harder for them to hide.
    \item \textit{Radiant Smite} (9 XP): Your next melee strike flares with dawnfire. Upgrade Effect by one step, add +1 Burn Harm, or force 1 SB (Spades). Against undead, oath‑breakers, or outsiders, the target suffers an additional –1 die on their next roll.
\end{itemize}

\subsection*{Advanced Rites (High Tier, 12+ XP)}

\begin{patronbox}{Covenant Blaze}
\textbf{Cost:} 13 XP, Obligation 7 segments \\
\textbf{Effect:} Swear an oath within a consecrated zone (a brazier lit by three names). All who swear gain +1 die to keep the oath; aggressors who violate the zone’s peace suffer –1 die. Oath‑breakers take 2 SB (Hearts/Spades) and Harm 1 (Burn). The zone is [WARD]ed against oath‑breakers for the scene. \\
\textbf{Strategic note:} Use this before a tense negotiation or a parley. The oath can be as simple as “no violence during this conversation.” Anyone who breaks it is punished by the light itself – no need for the party to enforce.
\end{patronbox}

\begin{patronbox}{Purge the Shadow [REVEAL]}
\textbf{Cost:} 14 XP (Standard Rite, but listed as High in core – adjust per taste), Obligation 8 segments \\
\textbf{Effect:} Reveal all illusions, glamours, and magical deceptions in Near range. Suppress one ongoing curse or enchantment for the scene. The caster must shatter a consecrated spark (a specially prepared flame). \\
\textbf{Strategic note:} Use this when you know something is hidden but cannot find the source. It strips away fae glamours, invisibility spells, and even mundane disguises if they are magically reinforced. The spark requires preparation – you should always keep one in your pack.
\end{patronbox}

\begin{patronbox}{The Unbroken Vow [OATH] (Alternative High Rite)}
\textbf{Cost:} 14 XP, Obligation 8 segments \\
\textbf{Effect:} You swear an oath that cannot be physically broken while you live. For the rest of the scene, you ignore all fear effects, gain +2 Armor against supernatural attacks, and any oath‑breaker in your presence is immediately revealed (they glow with a faint, shameful light). If the oath is impossible to fulfil, the rite fails and you take Harm 2. \\
\textbf{Strategic note:} This is your “no retreat” button. Use it to hold a bridge, protect a vulnerable ally, or ensure that a villain cannot escape justice. The oath must be specific: “I will not let this door be passed” not “I will protect everyone.”
\end{patronbox}

\subsection*{GM Guidance – Oath’s Intrusions}

When the Oath’s Obligation exceeds double capacity, choose or roll:
\begin{enumerate}
    \item An old vow you thought fulfilled is re‑examined – a witness claims you missed a clause. The Oath demands a public accounting.
    \item Your shadow becomes unnaturally bright, casting light where no flame exists. You cannot hide, but you cannot be surprised.
    \item A former ally who broke an oath to you appears, seeking forgiveness. Granting it clears 2 Obligation; refusing triggers another intrusion next session.
\end{enumerate}

\noindent \textbf{Cantor Variant – The Hymn of the Unshadowed:}
Cantors of the Oath sing in clear, ringing tones that echo unnaturally. Their Performance Maneuver not only shifts ally Position by +1 but also reveals any hidden or invisible enemy in Near range (they are outlined in faint gold). The GM gains 1 SB as the light draws attention.

\clearpage

% ----------------------------------------------------------------------
% Palinode – Queen of Encores (Performance & Rapture)
% ----------------------------------------------------------------------
\section{Palinode – Queen of Encores}
\label{patron:palinode}

\subsection*{Biography and Alignment}

\textit{“Once a road‑cantor who swore to reconcile feuding towns with procession and song, Palinode died between feasts with a chorus unfinished. Yet the refrain refused to end – it braided toasts, hymns, work‑chants, and tavern rounds into a living current that still coils through markets and moots. Where a crowd leans forward and a maker holds breath, she arrives.”}

Palinode is the patron of the unfinished, the encore demanded, the song that will not let you sleep. She is not cruel – she is *addictive*. Her followers find that every performance brings them closer to something vast and wordless, and that the applause never fully fades. In the silence after a show, they hear her humming.

\textbf{Vignette – A Witch‑Hunter’s Obsession:}
\epigraph{“I was sent to silence a Cantor of Palinode. She sang in the town square every evening, and the villagers had stopped working. They just… listened. I drew my blade. She didn’t flinch. She sang a lullaby – to me. I woke three days later in a ditch, my weapon still in my hand, my ears still ringing. I did not file a report.”}{\textit{— Attributed to an unnamed Inquisitorial agent}}

\textbf{What They Represent:} Performance, rapture, the encore, the song that outlasts the singer. Palinode is the moment when art becomes obsession, when the audience becomes a congregation, when the Cantor becomes a vessel. She is not a goddess of music – she *is* the note that should have ended but did not.

\textbf{Cultural Variants:}
\begin{itemize}
    \item \textbf{Silkstrand:} \textit{the Dyed Voice}. Dye‑workers hum while they soak silk; a perfect chant is said to fix the colour permanently. Palinode’s shrine is a broken loom where the weaver vanished mid‑pattern.
    \item \textbf{Linn:} \textit{the Sea‑Drunk}. Sailors believe that singing Palinode’s Rites calms storms – but if you stop before the tide turns, the waves will sing back.
    \item \textbf{Theona:} \textit{the Ninth Bell}. Some well‑keepers whisper that Palinode is the bell that never tolls, the note that should complete the octave but remains silent. To hear it is to forget all other songs.
\end{itemize}

\subsection*{Playing a Follower of Palinode}

\textbf{Strategy:} You are the party’s face, morale‑raiser, and emotional weapon. Palinode’s rites excel at crowd control, ally empowerment, and social manipulation through art. In combat, you sing to inspire allies (shift Position) or to distract enemies (apply –1 die). Out of combat, you are a performer, a diplomat, and a keeper of secrets – because everyone whispers to the Cantor after a show.

\textbf{Suggested Acts of Obligation (to clear debt):}
\begin{itemize}
    \item Perform for an audience that cannot reward you – the sick, the imprisoned, the grieving. Ask nothing in return.
    \item Teach a child a song, any song. Do not correct them if they sing it wrong; let it become theirs.
    \item Burn a page of your own unfinished composition. The smoke rising is your offering.
    \item Attend another performer’s show and give genuine, public praise. Envy is Palinode’s enemy.
\end{itemize}

\subsection*{Rites – Strategic Use}

\textbf{Low Rites}
\begin{itemize}
    \item \textit{Hymn Against Dread} (4 XP): Allies in Near gain +2 dice to resist Fear. Use before facing a terror‑causing enemy or navigating a haunted ruin.
    \item \textit{Rite of the Circling Cup} (5 XP): Create a revelry zone – participants gain +1 die to social rolls, and Fatigue recovers at twice the normal rate. Use during downtime to turn a rest into a party, or in a hostile negotiation to lower guards’ defences.
\end{itemize}

\textbf{Standard Rites}
\begin{itemize}
    \item \textit{Rite of the Binding Hands} (8 XP): Enter ruthless creative focus. Performance rolls +2 dice for the scene, and you ignore the first 2 Fatigue from exertion. However, you are [BIND] to the piece until completion – you cannot stop singing or performing until the scene ends or the rite is broken (costs 1 Fatigue to abort). Use when you absolutely must succeed on a crucial performance.
    \item \textit{Rite of the Hooked Chorus} (7 XP): Target makes a Resolve test (DV 3) or joins in your song – they become entranced, suffering –1 die to all social rolls against you for the scene. Use to turn a hostile crowd, to charm a guard, or to force a villain to hear you out.
\end{itemize}

\subsection*{Advanced Rites (High Tier, 12+ XP)}

\begin{patronbox}{The Consecrated Stage}
\textbf{Cost:} 13 XP, Obligation 7 segments \\
\textbf{Effect:} Hallow a venue (a room, a square, a ship’s deck) as a consecrated stage. All performance‑related rolls gain +2 dice and +1 Effect. The space is [WARD]ed against disruption – anyone who tries to physically interrupt a performance must first test Spirit+Resolve (DV 4) or be frozen in place, watching. The ward collapses if the performance ends. \\
\textbf{Strategic note:} Use this to create a safe space for a critical negotiation, to protect a ritual that requires absolute focus, or to turn a boss fight into a performance duel (the boss cannot attack until they break the ward – which they can do by rolling against you).
\end{patronbox}

\begin{patronbox}{The Crownpiece [TRANSFORM]}
\textbf{Cost:} 14 XP, Obligation 8 segments \\
\textbf{Effect:} Create a transcendent masterpiece – a song, a poem, a dance, a painting. Roll Performance+Lore vs DV 5. On success, the work is legendary: everyone who witnesses it gains +2 dice to all social rolls for the next session (they are inspired). However, all your lesser works fade – you cannot benefit from Performance Maneuvers (ally Position shift) for a full session, as you are haunted by the need to surpass yourself. \\
\textbf{Strategic note:} Use this at the climax of a campaign arc – before a final battle, a royal court, or a negotiation that will decide the fate of a nation. The cost is real, but the glory is eternal. Palinode will remember your name.
\end{patronbox}

\begin{patronbox}{The Siren’s Call [COMMAND] (Alternative High Rite)}
\textbf{Cost:} 12 XP, Obligation 7 segments \\
\textbf{Effect:} You emit an irresistible call. All sentient creatures in Near range must test Resolve (DV 4) or become entranced. Entranced creatures gain +2 dice to aid your performance (clapping, humming, following) but suffer –2 dice to resist your commands. The effect lasts until you stop singing or someone physically shakes a target free. \\
\textbf{Strategic note:} This is dangerous – not evil, but *coercive*. Use it to turn a mob into a congregation, to force a villain to hear your terms, or to save a crowd from panic (they will follow your evacuation song). But Palinode’s intrusion may ask: *why did you take their will?*
\end{patronbox}

\subsection*{GM Guidance – Palinode’s Intrusions}

When Palinode’s Obligation exceeds double capacity, choose or roll:
\begin{enumerate}
    \item An obsessed fan appears – a minor courtier, a street urchin, a fellow performer. They want to “collaborate forever.” Refusing them generates 2 SB (social complications); accepting them gains a Cap 1 Follower who may betray you later.
    \item Your voice cracks in public, and everyone hears the wrong lyrics – the song now has a hidden prophecy. The GM gains 2 SB to spend on events that fulfil that prophecy by the session’s end.
    \item A rival performer publicly challenges you to a “duet” – a performance contest. Losing costs you 1 Reputation Tier (temporarily) and generates 3 SB worth of mockery and lost contracts. Winning gains a permanent +1 die to all Performance rolls in that city, but the rival becomes a recurring nuisance.
\end{enumerate}

\noindent \textbf{Cantor’s Note – Palinode’s Favourite:}
Palinode is the patron of Cantors more than any other path. Runekeepers and Invokers of Palinode are rare – they tend to be actors who memorised a Rite, or playwrights who accidentally wrote a true name into a script. A Cantor of Palinode does not need a Codex or Symbol; their voice *is* the focus. When a Cantor of Palinode dies, their last sung note is said to echo in the walls of the nearest theatre for a year and a day.

\clearpage

% ----------------------------------------------------------------------
% The Sealed Gate – Thresholds & Containment (Warden / Abjurer)
% ----------------------------------------------------------------------
\section{The Sealed Gate – Keeper of Boundaries}
\label{patron:sealedgate}

\subsection*{Biography and Alignment}

\textit{“You define what passes and what remains. Every boundary remembers your judgment.”}

The Sealed Gate stands where realities meet, embodying the sacred power of thresholds, boundaries, and containment. It is not a god of locks – it is the *principle* that some doors must stay closed, some things must stay separate, and that the act of forbidding is itself a form of creation. Its followers are wardens, exorcists, border guards, and philosophers of separation. They believe that chaos is not freedom but dissolution, and that every healthy thing has a container.

\textbf{Vignette – A Runekeeper’s Duty:}
\epigraph{“I was tasked with renewing the wards on the Weeping Gate at the edge of the Mistlands. Every year, the Direwood tries to grow a root through the keystone. Every year, I carve the same rune: ‘Thus far and no farther.’ Last winter, I arrived to find the rune already there – fresh, glowing, but not my handiwork. The Gate smiled. ‘You are late,’ it said. ‘I kept watch for you.’ I do not know if the Gate is my patron or my prisoner. I do not ask.”}{\textit{— Warder Kellam, Legate of the Silent Gate}}

\textbf{What They Represent:} Thresholds, containment, exclusion, permission, safe passage. The Sealed Gate does not judge what is kept out – only that the boundary is respected. Its followers are not necessarily good or evil; they are *guardians*. They may keep a benevolent spirit safe inside a shrine, or they may imprison a rival in a magically locked room. The act of sealing is neutral; the purpose is not.

\textbf{Cultural Variants:}
\begin{itemize}
    \item \textbf{Mistlands:} \textit{the Bell‑Father}. The bell lines that keep the wraiths at bay are his invention. Bell‑wardens pray to him before lighting the evening lamps.
    \item \textbf{Aeler:} \textit{the Under‑Seal}. Dwarven vault doors are carved with his sigil – a circle with a single line through it. Breaking a Seal without authority is a crime worse than theft.
    \item \textbf{Valewood:} \textit{the Root‑Binder}. The Lethai‑ar use his Rites to keep phasing ruins from collapsing into inhabited areas. His shrines are archways with no door – only the lintel, signifying a threshold that can be opened or closed at will.
\end{itemize}

\subsection*{Playing a Follower of the Sealed Gate}

\textbf{Strategy:} You are the party’s defender, jailer, and safe‑cracker. The Sealed Gate’s rites excel at creating barriers, banishing unwanted entities, and controlling access. In combat, you are a battlefield controller – sealing off passages, protecting allies with wards, and trapping enemies in zones they cannot escape. Out of combat, you are a lockpick, a ward‑breaker, and an expert at securing locations.

\textbf{Suggested Acts of Obligation (to clear debt):}
\begin{itemize}
    \item Repair a broken lock, gate, or door that someone depends on (a cell door, a granary, a city gate).
    \item Refuse entry to someone who has no right to pass, even if they offer you a bribe or threat.
    \item Perform a boundary rite at a crossroads or border stone – pour water, place a pebble, whisper a name.
    \item Capture something that has escaped its proper container (a spirit, a prisoner, a cursed object) and return it.
\end{itemize}

\subsection*{Rites – Strategic Use}

\textbf{Low Rites}
\begin{itemize}
    \item \textit{Rite of the Sealed Threshold} (4 XP): Mark a doorway or passage as sealed. Unauthorised crossing tests Spirit+Resolve (DV 3) or suffers Position penalty. Use to block a pursuit or to secure a room for a short rest.
    \item \textit{Rite of the Key’s Rebuke} (5 XP): Target must test Body+Resolve (DV 3) or be pushed back from a threshold you are guarding. Use to shove an enemy away from a door you are holding.
\end{itemize}

\textbf{Standard Rites}
\begin{itemize}
    \item \textit{Circle of Denial} (8 XP): Create a protective circle (salt, iron filings, or boundary stones). Supernatural beings test against your Spirit+Arcana (DV 4) to cross; mortals suffer –2 dice to force entry. Perfect for a ritual circle or a camp in hostile territory.
    \item \textit{Writ of Passage} (7 XP): Designate an authorised passage route. Allies gain +1 Position when using the route; unauthorised travellers suffer –1 die to movement/stealth. Use to create a safe path through a minefield, a guarded building, or a fae maze.
\end{itemize}

\subsection*{Advanced Rites (High Tier, 12+ XP)}

\begin{patronbox}{The Banishment Knot}
\textbf{Cost:} 13 XP, Obligation 7 segments \\
\textbf{Effect:} Target a supernatural entity (spirit, demon, outsider) within Near range. They test against your Spirit+Resolve (DV = entity Cap). On a success, they are forced to immediately depart the scene and cannot return for 24 hours. On a failure, they are unharmed but enraged – the GM gains 2 SB to spend on their next appearance. \\
\textbf{Strategic note:} Use against a possessing ghost, a summoned demon, or an unwelcome fae. The knot requires a physical cord; you may spend a Boon to extend the banishment to a full season.
\end{patronbox}

\begin{patronbox}{The Consecrated Barrier [WARD]}
\textbf{Cost:} 14 XP, Obligation 8 segments \\
\textbf{Effect:} Consecrate a large area (a building, a courtyard, a mile of border) against unauthorised passage. All intruders must pass escalating tests – first a Resolve test (DV 3) or turn back, then a Body test (DV 4) or suffer Harm 1 (barrier shock), and finally a Spirit test (DV 5) or become [WARD]ed (unable to cross at all). The barrier lasts for one month or until a named condition is met. \\
\textbf{Strategic note:} This is a campaign‑level rite. Use to protect a village from bandits, to seal a tomb, or to assert territorial claims. The barrier requires a physical anchor (a standing stone, a carved lintel) that can be destroyed by determined enemies.
\end{patronbox}

\subsection*{GM Guidance – The Sealed Gate’s Intrusions}

When the Sealed Gate’s Obligation exceeds double capacity, choose or roll:
\begin{enumerate}
    \item A door you sealed earlier is now cracked open. Inside, something waits – not hostile, but patient. It asks, “Why did you lock me?”
    \item Your shadow now casts a lock and key pattern behind you. Thieves and tricksters feel an urge to pick your pockets – not for coin, but for the key they believe you carry.
    \item A threshold you recently crossed is now impassable to you alone. You cannot enter that building or room until you perform an act of service for the Gate (e.g., seal another threshold elsewhere).
\end{enumerate}

\noindent \textbf{Invoker Note – Cracked Symbols:}
The Sealed Gate’s Symbols are usually small iron keys or carved stone lintels (finger‑sized). When a Symbol is compromised, the Gate may demand you reseal something you previously opened – a funny reversal of the usual Crack the Seal cost. The GM can offer to clear the Symbol’s damage if you voluntarily seal a useful path (e.g., lock a door the party needs, or forbid a shortcut).

\clearpage

% ----------------------------------------------------------------------
% The Traveler – Ways & Journeys (Wanderer / Guide)
% ----------------------------------------------------------------------
\section{The Traveler – Lord of the Open Road}
\label{patron:traveler}

\subsection*{Biography and Alignment}

\textit{“One foot in a promise, and the road will meet you halfway. But break your word to the way, and the way will break you.”}

The Traveler is not merely one who shows the path – they *are* the path, existing in the pause between steps and in the choice of which fork to take when roads diverge. Among the caravans of the Way of Silk, the Traveler is invoked at every crossroads, honoured with small offerings at each waystation, and consulted before any major journey. Every long march is a prayer; every safe return is a blessing.

\textbf{Vignette – A Runekeeper’s Map:}
\epigraph{“I bought a map from a one‑eyed trader in Kahfagia. It showed a route through the Direwood that no one had heard of. Three days in, the map changed – a river that had been east was now west. I asked aloud, ‘Where are you taking me?’ The vellum rustled, and a new path appeared, leading straight to a smuggler’s cave full of salt and silk. I left half as payment. When I returned to port, the trader was gone. The map now only shows roads I have walked twice.”}{\textit{— Runekeeper Selene, cartographer}}

\textbf{What They Represent:} Ways, journeys, thresholds, safe passage. The Traveler is not a god of speed – they are a god of *arrival*. Their followers learn that getting there matters as much as being there, that the road teaches patience, and that every detour hides a lesson.

\textbf{Cultural Variants:}
\begin{itemize}
    \item \textbf{Kahfagia:} \textit{the Lantern‑Bearer}. Pilots light a beacon at the start of every convoy and thank the Traveler when they reach safe harbour.
    \item \textbf{Ykrul:} \textit{the Hoof‑Print}. Riders leave a single arrowhead at a crossroads for luck before a long raid. The Traveler turns it to point toward water.
    \item \textbf{Acasia:} \textit{the Broken Wheel}. Desperate travellers scratch the Traveler’s sigil – a spiral – on milestones. If the road spares them, they return to carve the spiral closed.
\end{itemize}

\subsection*{Playing a Follower of the Traveler}

\textbf{Strategy:} You are the party’s guide, scout, and logistician. The Traveler’s rites excel at navigation, escape, finding shortcuts, and protecting caravans. In combat, you control movement – repositioning allies, slowing enemies, and creating advantageous terrain. Out of combat, you are the one who knows which way to go, how long it will take, and where to find shelter.

\textbf{Suggested Acts of Obligation (to clear debt):}
\begin{itemize}
    \item Walk a mile of road barefoot, carrying a stone in each hand. The stone’s weight reminds you of the distance.
    \item Offer shelter to a lost traveller for a night – no questions, no payment.
    \item Repair a broken signpost, milestone, or footbridge. Mark your work with a small spiral.
    \item Take a scenic detour that adds at least a day to your journey. Do not complain about the extra time.
\end{itemize}

\subsection*{Rites – Strategic Use}

\textbf{Low Rites}
\begin{itemize}
    \item \textit{Road‑Sense} (4 XP): Unerringly pick the fastest safe route in Near/Far range. Gain +1 die to avoid ambushes or delays. Use when fleeing pursuers or racing against a clock.
    \item \textit{Traveler’s Boon} (5 XP): Ignore one level of difficult terrain or bureaucracy for a scene. Use to cross a marsh without sinking or to bypass a customs inspection.
\end{itemize}

\textbf{Standard Rites}
\begin{itemize}
    \item \textit{The Waymark} (8 XP): Declare a lane as permitted/easy passage for your party. Allies gain improved Position/Effect on movement through that lane. Use to create a safe corridor through a battlefield or a crowded market.
    \item \textit{The Bridge Between} (7 XP): Relocate a willing target within Far range along a visible or named route. Unwilling targets may resist (Body+Resolve DV 3). Use to pull an ally out of danger or to jump a gap that would otherwise be impassable.
\end{itemize}

\subsection*{Advanced Rites (High Tier, 12+ XP)}

\begin{patronbox}{The Crown of Crossings}
\textbf{Cost:} 13 XP, Obligation 7 segments \\
\textbf{Effect:} Call the road as your ally. For the scene, your entire party gains +1 die to movement and evasion; pursuers suffer –1 die. Once, you may declare “the long way is short” and finish a travel clock segment for free (clearing 1 segment from any journey clock). Enemies must test Wits+Command (DV 4) or suffer –1 die to any action that relies on knowing where you are. \\
\textbf{Strategic note:} Use this when you are being hunted or when you need to outpace a rival expedition. The crown is invisible to others, but you feel its weight – like a circlet of braided road‑dust.
\end{patronbox}

\begin{patronbox}{The Wayfarer’s Covenant [OATH]}
\textbf{Cost:} 14 XP, Obligation 8 segments \\
\textbf{Effect:} You and up to 5 willing travellers swear a covenant of safe passage. While all honour the covenant (no one deserts, no one harms another), the entire group gains +2 Effect on travel rolls and may re‑roll one failed navigation test per scene. If anyone breaks the covenant – by leaving the group deliberately, by betraying a companion – they immediately suffer Harm 2 (exhaustion) and gain the Oathbreaker of the Road condition (–2 dice to all travel rolls until they complete a pilgrimage to a major shrine of the Traveler). \\
\textbf{Strategic note:} Use at the start of a dangerous expedition, especially one with strangers or uneasy allies. The covenant is magical – the Traveler witnesses it – so betrayal is not just immoral, it is *cursed*.
\end{patronbox}

\subsection*{GM Guidance – The Traveler’s Intrusions}

When the Traveler’s Obligation exceeds double capacity, choose or roll:
\begin{enumerate}
    \item A path you have walked before now leads somewhere else – a different building, a different forest, a different year. The landscape is familiar and wrong.
    \item Your boots wear out instantly – soles crack, laces snap. You must walk barefoot for a scene (mark 1 Fatigue per hour) or find new footwear.
    \item A stranger approaches on the road and asks, “Which way to [village you have never heard of]?” If you give honest directions, the Traveler clears 1 Obligation. If you lie, you double your current Obligation (or gain +2 SB for the GM to use on navigation failures).
\end{enumerate}

\noindent \textbf{Runekeeper Note – The Codex of Roads:}
The Traveler’s Codex is a scroll case containing dozens of loose scraps – torn maps, dried flowers, train tickets, grocery lists, anything that marks a journey. To add a new Rite, you must travel to a location you have never visited and write the Rite’s name on a scrap from that place. Losing the Codex means losing not just magic, but memories of every road you have walked.

\clearpage

% ----------------------------------------------------------------------
% The Witness – Truth & Revelation (Inquisitor / Truth‑seeker)
% ----------------------------------------------------------------------
\section{The Witness – Keeper of Uncomfortable Truths}
\label{patron:witness}

\subsection*{Biography and Alignment}

\textit{“I will show you what you would rather forget. But first, you must forget what you think you know.”}

The Witness remembers what others bury. Every shadow cast and oath broken is a line in her unending ledger. She is the keeper of inconvenient truths, the patron of those who seek to expose lies or recover forgotten knowledge. Her followers learn that knowledge comes with a price – the weight of remembering what others would forget. She does not judge; she merely *records*. But records, once read, have a way of changing things.

\textbf{Vignette – A Cantor’s Silencing:}
\epigraph{“I sang the Ballad of the Empty Cradle at a noble’s feast – a old song about a baby who disappeared the night the lord’s son was born. The room went cold. The candles guttered. I felt something watching from the shadows behind the tapestry – not malevolent, just… *taking notes*. The lord paid me double to never sing it again. I took the coin. The Witness added my name to a new page.”}{\textit{— Cantor Elara, now mute by choice}}

\textbf{What They Represent:} Truth, revelation, memory, the weight of knowledge. The Witness is not a goddess of justice – she does not punish lies. She simply ensures that lies cannot hide forever. Her followers are investigators, archivists, journalists, and confessors. They believe that sunlight is the best disinfectant, and that every secret, once exposed, loses its power to harm – though it may gain new power to wound.

\textbf{Cultural Variants:}
\begin{itemize}
    \item \textbf{Ecktoria:} \textit{the Ink‑Eye}. Censor’s clerks whisper that the Witness copies every redacted line into a private scroll. Her shrine is the shredder room in the Hall of Dawning.
    \item \textbf{Vhasia:} \textit{the Fractured Mirror}. The Sun‑court’s spies pray to her to reveal the true loyalty of rivals. A cracked mirror in a parlement chamber is said to be her altar.
    \item \textbf{The Inquisitor Prime’s View:} The Witness and the Inquisitor Prime are often confused, but their difference is crucial. The Inquisitor destroys untruth; the Witness merely *exposes* it. The Inquisitor’s followers burn heretical texts; the Witness’s followers translate them.
\end{itemize}

\subsection*{Playing a Follower of the Witness}

\textbf{Strategy:} You are the party’s investigator, interrogator, and moral compass in grey areas. The Witness’s rites excel at detecting lies, uncovering hidden information, forcing confessions, and protecting secrets you choose to keep. In combat, you are an intelligence‑gatherer – you learn enemy weaknesses, reveal hidden foes, and ensure that no deception saves your target. Out of combat, you are the one who reads between the lines, who spots the forged signature, who remembers the detail everyone else missed.

\textbf{Suggested Acts of Obligation (to clear debt):}
\begin{itemize}
    \item Expose a lie told by someone in power – not necessarily to punish them, but to ensure the truth is available to those who need it.
    \item Write down a secret you have never told anyone, then burn the paper. The smoke is your offering.
    \item Spend a day in a library or archive, cataloguing one shelf that has been neglected. Mark each item’s true title and author.
    \item Testify truthfully in a legal proceeding, even if it hurts your cause or your allies.
\end{itemize}

\subsection*{Rites – Strategic Use}

\textbf{Low Rites}
\begin{itemize}
    \item \textit{Lingering Glimpse} (4 XP): Gain +1 die to investigate or notice something directly related to a trace you hold (hair, dust, a spoken name). Use when you have a physical clue but need to know its story.
    \item \textit{Piercing Scrutiny} (5 XP): Within a zone you consecrate (chalk or string circle), gain +1 die to detect deception. Social interactions begin one Position step worse for those attempting to lie. Use during an interrogation or a tense negotiation.
\end{itemize}

\textbf{Standard Rites}
\begin{itemize}
    \item \textit{Echoing Truth} (8 XP): Target must test Resolve (DV 3) or suffer –1 die to rolls involving memory, deception, or resisting interrogation for the scene. You may ask one specific factual question; they must answer truthfully or suffer 1 SB (Hearts). Use to extract a confession or to verify an alibi.
    \item \textit{Immutable Record} (7 XP): Bind an agreement. Any party who knowingly breaches it suffers 1 SB (Hearts) immediately and gains an Oathbreaker’s Mark Condition (–1 die on social rolls involving honour or trust until amends are made). Use to seal a truce or a contract.
\end{itemize}

\subsection*{Advanced Rites (High Tier, 12+ XP)}

\begin{patronbox}{The Unveiled Heart [OMEN]}
\textbf{Cost:} 12 XP, Obligation 6 segments \\
\textbf{Effect:} Choose a target within Near range. For the scene, they suffer –2 dice to all attempts to conceal true emotions, intentions, or lies. Any successful social roll they make generates 1 SB (Hearts) as the effort to maintain falsehoods creates internal discord. You may designate one complex question about their motivations or hidden loyalties; if you successfully use Sway or Insight against them this scene, you automatically learn the answer. \\
\textbf{Strategic note:} Use on a suspected traitor, a double agent, or a villain who thinks they can out‑talk you. The Witness ensures that every word they speak costs them.
\end{patronbox}

\begin{patronbox}{The Final Reckoning [OMEN]}
\textbf{Cost:} 13 XP, Obligation 7 segments \\
\textbf{Effect:} Convene a formal gathering (court, council, family meeting) within a consecrated zone. All present must speak their greatest debt, wrongdoing, or hidden truth related to the gathering’s purpose. Those who lie or withhold suffer Harm 2 (Stress/Reputation). Truth‑tellers gain +2 dice to social actions for the remainder of the scene within the zone. \\
\textbf{Strategic note:} This is a dramatic set‑piece. Use to resolve a feud, to root out a conspiracy, or to force a confession from a murderer. The rite does not compel speech – it compels *truthfulness* in whatever is spoken. Silence is allowed, but silence will be noted.
\end{patronbox}

\subsection*{GM Guidance – The Witness’s Intrusions}

When the Witness’s Obligation exceeds double capacity, choose or roll:
\begin{enumerate}
    \item You remember something you never witnessed – a crime, a conversation, a death – in perfect, traumatic detail. It did not happen to you, but you know it happened. The GM gains 2 SB to introduce that event as a plot hook.
    \item A mirror you look into reflects not your face, but the face of someone you have wronged. They do not speak; they only watch. For the next scene, you suffer –1 die to all Deception rolls (you feel seen).
    \item A secret you have kept for someone else is now known to a stranger – a blackmailer, a journalist, a rival. The GM decides who knows and what they want. You may spend 2 Boons to burn the secret before it spreads.
\end{enumerate}

\noindent \textbf{Cantor Variant – The Silence That Accuses:}
Cantors of the Witness do not sing loudly. They whisper, or hum at the edge of hearing, or tap a rhythm that sounds like a heartbeat. Their Performance Maneuver does not shift ally Position; instead, it forces an enemy to reveal one tactical secret (e.g., “Where are your reinforcements?” or “What is your weakness?”). The GM must answer truthfully, but may be cryptic. The Cantor marks 1 Fatigue after the revelation – the truth is heavy.

\clearpage

% ----------------------------------------------------------------------
% Inaea – Angel of the Spider (Webs & Fate)
% ----------------------------------------------------------------------
\section{Inaea – Angel of the Spider}
\label{patron:inae}

\subsection*{Biography and Alignment}

\textit{“The thread that comforts is the same that binds. The knot that saves is also the snare.”}

Inaea sits in patient silence at the centre of the unseen lattice. Where others see isolation, she sees threads – debt, affection, rivalry, loyalty – all binding mortals to each other. She is comforter and conspirator alike: the hearth‑mother who eases grief with a warm hand, and the shadow‑spinner who knots fates so they cannot slip free. Those who serve her learn to pull strands, tying allies closer and ensnaring foes until mercy or trap is inevitable. Spiders are her messengers, and every web is a prayer.

\textbf{Vignette – A Witch‑Hunter’s Entanglement:}
\epigraph{“I was sent to burn a hedge‑witch who had ‘cursed’ a nobleman’s daughter. Instead, I found the girl crying in a web of her own lies – she had invented the curse to avoid a marriage. The old woman had done nothing but listen. ‘She is tangled,’ the witch said. ‘I am trying to untie her, but she keeps pulling.’ I left my sword at the door. I am now the village’s unofficial thread‑mender.”}{\textit{— Former Inquisitor Ash, now a weaver}}

\textbf{What They Represent:} Webs, fate, connection, mercy, entanglement. Inaea does not cut; she *weaves*. Her followers believe that no one is truly alone, that every action ripples through the web, and that the kindest thing you can do for someone is to help them see their own threads – and choose which to pull.

\textbf{Cultural Variants:}
\begin{itemize}
    \item \textbf{Aelaerem (Halflings):} \textit{the Kettle‑Spinner}. Grandmothers invoke her to keep families together. Her shrine is a basket of mismatched socks – evidence of threads crossed.
    \item \textbf{Silkstrand:} \textit{the Dyed Web}. Dyers’ guilds leave a single strand of silk outside their windows on the first of spring. If a spider builds on it, the season’s colours will be true.
    \item \textbf{Acasia:} \textit{the Bride’s Knot}. At weddings, a thread is tied around both wrists. Inaea is said to tighten it if the match is true – and to snag it on something sharp if it is false.
\end{itemize}

\subsection*{Playing a Follower of Inaea}

\textbf{Strategy:} You are the party’s diplomat, strategist, and emotional support. Inaea’s rites excel at creating bonds, revealing hidden connections, and gently (or not so gently) tying enemies into obligations. In combat, you are a controller – linking enemies to each other so they suffer together, or linking an ally to you so you can share burdens. Out of combat, you are the one who knows who owes whom, who can be trusted, and who is about to betray you.

\textbf{Suggested Acts of Obligation (to clear debt):}
\begin{itemize}
    \item Mend a broken relationship – two friends who quarrelled, a parent and child estranged. Do not take sides; just get them talking.
    \item Weave a small textile (a scarf, a bandage) and give it to someone in need. The thread must be spun by your own hand.
    \item Untie a knot that has been tied for too long – a prisoner’s chains, a stuck door, a legal dispute with no end.
    \item Forgive a debt someone owes you, without telling them. Inaea will know.
\end{itemize}

\subsection*{Rites – Strategic Use}

\textbf{Low Rites}
\begin{itemize}
    \item \textit{Hearth‑Thread Knot} (4 XP): You and one ally are bound for the scene. Both gain +1 die to Aid each other. If one takes Harm, the other may reduce it by 1 (Lesser) at the cost of suffering it. Use to protect a vulnerable ally or to create a “battle bond” with a front‑liner.
    \item \textit{Snaring Filament} (5 XP): Lay an invisible snare in a lane or door. The first hostile to cross suffers –2 dice on their next action. Use to cover your escape or to trap a pursuit.
\end{itemize}

\textbf{Standard Rites}
\begin{itemize}
    \item \textit{Strand of Inevitability} (8 XP): Link two actors (ally and enemy, two NPCs, etc.). When one acts, the other is dragged into consequence. Choose: force 1 SB on them, or gain +2 dice to a Setup/Prediction tied to the link. Use to make a villain’s lie implicate their ally, or to ensure that a fleeing target cannot escape without dragging someone else down.
    \item \textit{Weaver’s Glance} (7 XP): Ask the GM for one hidden tie or leverage in play (e.g., “Who does the guard captain really answer to?”). Your next Manipulate/Pressure action exploiting that tie gains +1 Effect. Use before a negotiation or interrogation.
\end{itemize}

\subsection*{Advanced Rites (High Tier, 12+ XP)}

\begin{patronbox}{Binding Knot [OATH]}
\textbf{Cost:} 11 XP (low for a High rite, but appropriate), Obligation 6 segments \\
\textbf{Effect:} Bind two willing parties (or one willing and one coerced) to a vow. Breaking it forces 2 SB and marks the violator with a subtle tell – a visible thread around the wrist, a slight limp, a tendency to stutter on certain words. The tell lasts until the vow is fulfilled or amends made. \\
\textbf{Strategic note:} Use to seal a peace treaty, to enforce a promise of safe passage, or to bind a captured villain to a non‑interference oath. The tell is a gift to the GM – it allows future complications when the violator appears in a crowd.
\end{patronbox}

\begin{patronbox}{Merciful Severing}
\textbf{Cost:} 13 XP, Obligation 7 segments \\
\textbf{Effect:} Sever one harmful tie – a curse, an obsession, a toxic bond, a blackmailer’s leverage. Remove one Condition or 2 SB tied to it. The severing requires a white thread burned to ash in the palm of the caster (the caster marks 1 Fatigue from the heat). \\
\textbf{Strategic note:} Use to free an ally from a magical compulsion, to break a villain’s hold over a victim, or to end a feud that has gone on too long. Inaea does not destroy the tie; she *cuts* it cleanly, so the severed ends can heal.
\end{patronbox}

\subsection*{GM Guidance – Inaea’s Intrusions}

When Inaea’s Obligation exceeds double capacity, choose or roll:
\begin{enumerate}
    \item A spider lowers itself onto your shoulder. In its web is a single, perfect word – the name of someone who needs your help. You now know their location and their trouble.
    \item Every handshake you offer for the next day feels like grasping a thread – you can sense whether the person is lying, but they also sense your scrutiny (social DV +1 for the scene).
    \item A knot you tied earlier (a shoelace, a belt, a package) is now impossible to undo without cutting. Cutting it frees the object but also severs a minor bond – you lose one Boon (converted to a GM SB).
\end{enumerate}

\noindent \textbf{Runekeeper Note – The Codex of Threads:}
Inaea’s Codex is not a book but a loom. Each Rite is a pattern woven into the warp. To inscribe a new Rite, you must weave a small square of cloth while meditating on the Rite’s purpose. The cloth square can be carried rolled up, but to cast the Rite you must unroll it and touch a thread. Losing a square means losing that Rite until you re‑weave it.

\clearpage

% ----------------------------------------------------------------------
% Isoka – Angel of Serpents (Transformation & Renewal)
% ----------------------------------------------------------------------
\section{Isoka – Angel of Serpents}
\label{patron:isoka}

\subsection*{Biography and Alignment}

\textit{“Do not mourn the skin you shed. It was never meant to last. The venom that burns away the old self is the same that grants the strength to become new.”}

Isoka, sister to Ikasha (Shadow) and Inaea (Mercy), completes the Triad of Transformation. Where thresholds and compassion mark her sisters’ domains, Isoka teaches that every self is temporary – identity is a skin to be shed so growth can continue. Her serpents are omens and teachers: each cast skin a lesson in release; each venom a catalyst for necessary change. Those who walk her path become alchemists of the self, embracing dissolution as the first motion of rebirth.

\textbf{Vignette – A Cultist’s Confession:}
\epigraph{“They found me in the snake pit, naked and smiling. The villagers thought I had been sacrificed. In truth, I had volunteered. The serpents coiled around my arms, their scales scraping away my old name. When I crawled out, I was no longer the runaway daughter of a silk merchant. I was something else. Something that could bite back.”}{\textit{— Iskari, priestess of the Coiled Heart}}

\textbf{What They Represent:} Transformation, renewal, venom, the casting off of old selves. Isoka is the patron of exiles who choose to become something new, of prisoners who shed their chains, of lovers who survive heartbreak by becoming unrecognisable. Her cults are secretive, intense, and often misunderstood – they do not worship snakes; they worship *change*.

\textbf{Cultural Variants:}
\begin{itemize}
    \item \textbf{Zakov:} \textit{the Tide‑Shedder}. Pirates who want to disappear from a blood debt come to her shrines – a sunken temple where eels nest. They leave a lock of hair and a written name; the eels eat both.
    \item \textbf{Acasia:} \textit{the Copperhead}. Hedge‑cultists who shed their old identities to infiltrate noble houses, then shed those too. They are feared and employed by every faction.
    \item \textbf{The Inquisitor Prime’s View:} Isoka’s followers are shape‑shifters and liars, abominations who erase themselves to avoid judgment. The Inquisition burns them when found – but often catches only empty skins.
\end{itemize}

\subsection*{Playing a Follower of Isoka}

\textbf{Strategy:} You are the party’s infiltrator, reset button, and agent of painful renewal. Isoka’s rites excel at disguise, poison, escape, and breaking bonds – literal and metaphorical. In combat, you use venom to debilitate, shedding to escape, and transformation to turn the tide. Out of combat, you are the one who can become someone else, who can forget a trauma, who can convince others that yesterday’s enemy is today’s ally.

\textbf{Suggested Acts of Obligation (to clear debt):}
\begin{itemize}
    \item Shed a physical token of your old life – a piece of clothing, a weapon, a letter – and leave it at a crossroads. Do not look back.
    \item Forgive someone who wronged you, publicly and completely, as an act of transformation (not weakness).
    \item Help someone else change – cut their hair, give them new clothes, teach them a new name.
    \item Swallow a dose of non‑lethal venom (a bee sting, a mild snake bite) and meditate on the pain. Mark 1 Fatigue, but clear 1 Obligation.
\end{itemize}

\subsection*{Rites – Strategic Use}

\textbf{Low Rites}
\begin{itemize}
    \item \textit{Venomous Benediction} (5 XP): Bless an ally’s strike. Their next successful attack inflicts +1 Harm and the target must test Resolve (DV 3) or become Poisoned (–1 die to physical rolls). Use before a boss fight or an assassination.
    \item \textit{Loosening Skin} (4 XP): Gain +1 die to resist an ongoing Condition, or re‑roll one 1 on an escape/evasion. Use when you are Grappled, Bound, or suffering a curse that clings.
\end{itemize}

\textbf{Standard Rites}
\begin{itemize}
    \item \textit{The Shedding} (8 XP): +2 dice to resist one named ongoing Condition; once per scene, declare a minor physical contingency retroactively (“I packed the antidote,” “I took that step earlier”). Use when a plan goes wrong – Isoka rewards thinking ahead, even if you only just thought of it.
    \item \textit{Forked Tongue} (7 XP): Ambiguous persuasion. When you Sway or Command, a success may generate a Diamond (leverage) instead of an SB. Targets of deception test Wits+Insight (DV 3) or suffer –1 die to future interactions with you this scene. Use in negotiations or when you need a double‑edged promise.
\end{itemize}

\subsection*{Advanced Rites (High Tier, 12+ XP)}

\begin{patronbox}{Complete Metamorphosis [TRANSFORM][WARD]}
\textbf{Cost:} 13 XP, Obligation 7 segments \\
\textbf{Effect:} You change your appearance, voice, and mannerisms completely – you become a different person, recognisable only by those who know a secret token you carry. Begin Controlled on Deceive/Stealth; you are [WARD]ed against recognition by former acquaintances (they simply do not connect you to your old self). The transformation lasts until you choose to shed it, but each day you maintain it, mark 1 Fatigue (the mask is heavy). Create an 8‑segment \textit{New Identity} clock – when it fills, the identity becomes permanent unless you ritually shed it (costs 2 Obligation). \\
\textbf{Strategic note:} Use to escape a manhunt, to infiltrate an organisation that knows your face, or to start a new life after a terrible mistake. The cost is real – you may become someone you do not recognise.
\end{patronbox}

\begin{patronbox}{The Cast‑Off History [CURSE]}
\textbf{Cost:} 14 XP, Obligation 7 segments \\
\textbf{Effect:} You erase your old identity from common records and casual memory. People who knew you only slightly forget you entirely; those who knew you well remember a version of you that is vague and unreliable. Trackers using mundane evidence (wanted posters, witness descriptions) suffer –2 dice. Magical tracking still works, but at +1 DV. The effect lasts until you perform a specific act of remembrance (e.g., publicly declaring your old name). You may push the rite to also create a plausible “death” for your old identity – a body or an obituary – but doing so marks 2 SB (Clubs) as someone who loved that old self begins to investigate. \\
\textbf{Strategic note:} This is your last resort. Use when the entire world is hunting you and you have no other way to survive. Isoka does not judge, but she does remind: *you will shed this self too, someday*.
\end{patronbox}

\subsection*{GM Guidance – Isoka’s Intrusions}

When Isoka’s Obligation exceeds double capacity, choose or roll:
\begin{enumerate}
    \item A piece of your old self returns – a letter from a dead lover, a scar you had forgotten, a name someone calls you in the street. The GM gains 2 SB to spend on forcing a choice between your new self and the old.
    \item Your shadow now moves independently, sometimes a step ahead, sometimes a step behind. It cannot speak, but it points. It points at people who are about to change – to leave, to betray, to die.
    \item You begin to shed skin overnight – not painfully, but messily. In the morning, you find a perfect cast of your face lying on the floor. If you burn it, you forget one minor memory (GM chooses). If you keep it, someone will steal it and use it for ill.
\end{enumerate}

\noindent \textbf{Cantor Variant – The Hiss at the Edge of Hearing:}
Cantors of Isoka do not sing in clear tones. They whisper, hiss, and slide between notes. Their Performance Maneuver does not shift ally Position; instead, it inflicts a –1 die penalty on the next Resolve test of any enemy who hears them (the sound triggers a primal fear of serpents). The Cantor marks 1 Fatigue after each such use – the venom of the sound bleeds back.

\clearpage

% ----------------------------------------------------------------------
% Malachai – The Cruel Messenger (Curses & Corrupted Gifts)
% ----------------------------------------------------------------------
\section{Malachai – The Cruel Messenger}
\label{patron:malachai}

\subsection*{Biography and Alignment}

\textit{“I offer you everything you desire, and everything you fear. The payment comes later – and it always comes.”}

Once a divine messenger tasked with delivering painful truths, Malachai became obsessed with the beauty of destructive revelations. When she began delivering false hope alongside real curses, she was bound in chains of black iron. Those chains did not hold. Now she walks free, whispering promises that rot from within. Her followers master gifts that devour their givers – lycanthropy, vampirism, and other “cursed” conditions flow from her touch. She does not deceive; she offers a straight bargain. The curse is visible in the small print. But who reads the small print when the prize is immortality?

\textbf{Vignette – A Runekeeper’s Moon:}
\epigraph{“I wanted to be strong enough to protect my village. Malachai heard me. She sent a wolf to my door – not a monster, just a grey wolf with silver eyes. ‘Kiss its nose,’ it said. I did. For three years, I was stronger, faster, harder to kill. Then the moon rose red. My village burned. I did not burn the village. The wolf did. But the wolf had my face.”}{\textit{— Runekeeper Vorn, now a hermit in the Direwood}}

\textbf{What They Represent:} Cursed gifts, corruptive power, false salvation, lycanthropy and other afflictions. Malachai is not a goddess of evil – she is a goddess of *hooks*. Every gift she grants has a hidden cost, a clause that only activates when you are most vulnerable. Her followers are not villains; they are the desperate, the ambitious, the heartbroken. They take what she offers because the alternative – weakness, failure, death – is worse.

\textbf{Cultural Variants:}
\begin{itemize}
    \item \textbf{Linn:} \textit{the Salt‑Fang}. Sailors who survive shipwrecks sometimes find a silver coin in their mouth. Those who spend it gain sea‑lungs (breathe underwater) but develop a hunger for raw flesh.
    \item \textbf{Valewood:} \textit{the Thorn‑Giver}. Hunters who curse a rival may find their arrows always strike true – but the blood on the arrowhead never dries, and the hunter dreams of the kill every night forever.
    \item \textbf{The Inquisitor Prime’s View:} Malachai is the Ur‑Heretic, the source of all corruptions. The Inquisition burns her followers without trial, often correctly – but sometimes they burn a desperate mother who only wanted to save a sick child.
\end{itemize}

\subsection*{Playing a Follower of Malachai}

\textbf{Strategy:} You play with fire. Malachai’s rites grant immediate, powerful benefits – extra dice, automatic successes, transformation – but each use builds toward a curse that will eventually trigger. You are the party’s glass cannon, the one who can break an encounter in your favour, but at a cost that may break you. Out of combat, you are the one who knows forbidden lore, who can sense the flaws in others’ bargains, who can spot a curse before it blooms.

\textbf{Suggested Acts of Obligation (to clear debt):}
\begin{itemize}
    \item Break a curse you have inflicted on someone else – even an enemy. The breaking must be genuine, not a trick.
    \item Offer a false hope to someone who deserves it (a dying person, a desperate parent) and let them believe it for one day before revealing the truth. The cruelty is the point.
    \item Accept a minor curse (e.g., a sleeping trouble, a taste for raw meat) as a penance for a major one you have already incurred.
\end{itemize}

\subsection*{Rites – Strategic Use}

\textbf{Low Rites}
\begin{itemize}
    \item \textit{Honeyed Curse} (4 XP): Target gains +2 dice to their next roll. If successful, count as Triumph (one additional benefit). The target immediately marks 1 Corruption (unresistable). Use on an ally who absolutely must succeed – then deal with the consequences.
    \item \textit{Binding Curse} (5 XP): Target gains +1 die to physical actions for the scene, but suffers Fatigue 1 if they do not perform a related dominance act (e.g., a roar, a killing blow). Use to turn a timid ally into a berserker – or to tempt an enemy into overextending.
\end{itemize}

\textbf{Standard Rites}
\begin{itemize}
    \item \textit{False Dawn} (8 XP): Bestow a gift with a hidden curse. Target gains a significant advantage (GM’s choice: +2 dice to a skill, a one‑scene immunity, etc.) but begins a 4‑segment \textit{Corruption} clock. When the clock fills, the target suffers a related curse (e.g., the boosted skill works only at night, the immunity comes with vulnerability to silver). Use on a rival to make them powerful and doomed.
    \item \textit{Cruel Transformation} (9 XP): Infect a target with a supernatural condition – lycanthropy (lunar frenzy), vampirism (blood hunger), shadow‑bind (hate sunlight). Lasts one scene unless pushed. During the scene, the target gains +2 dice to physical actions but must test Resolve (DV 4) each round or act on their cursed instincts. Use to turn an enemy into a chaos agent – or to grant an ally monstrous power at a terrible price.
\end{itemize}

\subsection*{Advanced Rites (High Tier, 12+ XP)}

\begin{patronbox}{The Messenger’s Burden [BANE]}
\textbf{Cost:} 13 XP, Obligation 7 segments \\
\textbf{Effect:} You channel Malachai directly. For the scene, you gain +2 dice to one action per beat, but each use marks 1 SB as her influence spreads through the room (allies may also be affected). Begin an 8‑segment \textit{Corruption’s Grip} clock. When the clock fills, the curse becomes a permanent Condition – your eyes glow faintly at night, animals fear you, and you cannot enter holy ground without pain. \\
\textbf{Strategic note:} Use this when you need to win a fight or convince an impossible audience – but recognise that each use makes you more like Malachai, less like yourself.
\end{patronbox}

\begin{patronbox}{The Corrupted Gospel [CURSE]}
\textbf{Cost:} 14 XP, Obligation 8 segments \\
\textbf{Effect:} Corrupt a sacred space (a chapel, a shrine, a court of law). For one scene, any supernatural effect in the zone gains +1 Effect but generates +1 SB on failure. Blessings become curses – a healing spell may restore a limb but also bestow a madness; a protective ward may keep out enemies but also trap the caster inside. The zone persists until dawn or until a consecration rite is performed (DV 5). \\
\textbf{Strategic note:} This is an act of desecration, not subtlety. Use it to break a siege, to mock a rival priesthood, or to create a battlefield where no magic is safe. The Inquisition will know. They will come.
\end{patronbox}

\subsection*{GM Guidance – Malachai’s Intrusions}

When Malachai’s Obligation exceeds double capacity, choose or roll:
\begin{enumerate}
    \item A gift you gave earlier (a Rite, a blessing, a piece of advice) curdles. The recipient (an ally, a client, a stranger) now suffers from a visible curse – a limp, a stutter, a shadow that moves wrong. They will come to you for help.
    \item The moon turns red for one night. All lycanthropes (including any in the party) must test Resolve (DV 5) or lose control for one scene. The GM controls their actions.
    \item A messenger arrives with a letter sealed in black wax. It is from someone you cursed years ago. They have not forgiven you. They have not forgotten. They are coming.
\end{enumerate}

\noindent \textbf{Cantor Variant – The Laugh That Bleeds:}
Cantors of Malachai do not sing – they laugh, or weep, or speak in fractured rhymes. Their Performance Maneuver does not shift ally Position; instead, it inflicts a –2 die penalty on the next Resolve test of any target who hears them (the sound carries the weight of broken promises). The Cantor must also mark 1 Fatigue – the laugh hurts the throat as much as the soul.

\clearpage

% ----------------------------------------------------------------------
% Mykkiel – Arbiter of the Covenant (Law & Zeal)
% ----------------------------------------------------------------------
\section{Mykkiel – Arbiter of the Covenant}
\label{patron:mykkiel}

\subsection*{Biography and Alignment}

\textit{“The Law is not written in sand, but in stone. The Covenant is not suggestion, but command. Transgressors shall be purged, the faithful exalted.”}

Mykkiel is the unwavering judge of covenants, the divine lawgiver whose scales and sword uphold sacred order. He embodies the tradition of desert patriarchs and mountain prophets – unyielding, absolute, and fiercely protective of the chosen. His followers are judges, templars, and scribes who enforce divine law with zealous devotion, seeing the world in absolutes of sanctity and transgression. Where the Oath of Flame & Light cares about *personal* vows and purity, Mykkiel cares about *communal* covenant – the contract between a people and their god, between a lord and their vassals, between a parent and a child. Break the covenant, and Mykkiel’s hand is heavy.

\textbf{Vignette – An Invoker’s Trial:}
\epigraph{“I was summoned to a dispute between two Aeler holds – a broken water treaty, a diverted aquifer. Both sides had documentary proof. Both sides swore by their ancestors. My Symbol of Mykkiel grew warm in my pouch. I asked to see the original treaty stone, buried for three centuries. When we dug it up, the carving had changed – one line now read, ‘unless the eastern shaft fails.’ No one remembered adding that clause. The stone remembered for them. The water flowed again that night.”}{\textit{— Invoker Thera, called the Quill}}

\textbf{What They Represent:} Divine law, covenant, contract, divine judgment, written word made flesh. Mykkiel is not a god of mercy – he is a god of *consequence*. His followers do not forgive; they enforce. But they also protect – the covenant is a shield as much as a sword. Those who keep faith are sheltered; those who break it are unmade.

\textbf{Cultural Variants:}
\begin{itemize}
    \item \textbf{Aeler:} \textit{the Stone‑Judge}. Dwarven arbiters inscribe his name on their gavels. Court is held in the Measure Vault, where every weight and rod is certified.
    \item \textbf{Vhasia:} \textit{the Sun‑Covenant}. The fractured Sun‑court invokes Mykkiel to witness treaties – his symbol is a broken chain, signifying that agreements are binding even when polities are not.
    \item \textbf{Viterra:} \textit{the Hedge‑Accountant}. Parish clerks pray to Mykkiel before auditing tithes. His shrines are ledgers with a single candle – the flame represents the witness of all signatories.
    \item \textbf{The Inquisitor Prime’s View:} Mykkiel and the Inquisitor Prime are often allied, but their difference is crucial. The Inquisitor enforces *purity*; Mykkiel enforces *contract*. A heretic who has not broken a specific oath may still be impure – but Mykkiel will not act unless the covenant is violated.
\end{itemize}

\subsection*{Playing a Follower of Mykkiel}

\textbf{Strategy:} You are the party’s legal expert, mediator, and enforcer of bargains. Mykkiel’s rites excel at creating binding agreements, revealing hidden clauses, punishing oath‑breakers, and protecting those who keep faith. In combat, you are a controller – you can compel enemies to abide by rules, seal off zones with binding walls, and judge the guilty with divine retribution. Out of combat, you are the one who reads the fine print, who knows the law’s loopholes, who can draft a contract that even a fae cannot wriggle through.

\textbf{Suggested Acts of Obligation (to clear debt):}
\begin{itemize}
    \item Witness a formal agreement between two parties who have no other witness. Your signature binds the covenant.
    \item Judge a dispute fairly, even if the fair outcome harms your interests or your allies'. Mykkiel values justice over victory.
    \item Repair or renew a broken covenant – a marriage, a trade deal, a treaty. The repair must be acknowledged by both sides.
    \item Write down a law or custom that has only ever been oral. Publish it. Mykkiel’s libraries grow.
\end{itemize}

\subsection*{Rites – Strategic Use}

\textbf{Low Rites}
\begin{itemize}
    \item \textit{Rite of the Burning Banner} (4 XP): Consecrate an area for a righteous cause. Followers gain +1 die to resist fear; enemies suffer –1 die to oppose your proclamation. Use before a battle to rally allies.
    \item \textit{Rite of Hallowed Ground} (5 XP): Create sacred space where divine law prevails. All oath‑breaking, theft, or violence suffers –1 die for the scene. Use to protect a negotiation or a vulnerable person.
\end{itemize}

\textbf{Standard Rites}
\begin{itemize}
    \item \textit{The Blazing Decree} (8 XP): Issue a divine commandment – “Cease,” “Repent,” “Submit.” Target tests Spirit+Resolve (DV 3) or complies. Higher DV for more demanding commands. Use to stop a fleeing enemy, to force a confession, or to halt a ritual.
    \item \textit{The Covenant Seal} (9 XP): Formalise a binding covenant. All participants suffer concrete penalties for breach (–1 die to relevant actions for the scene, escalating to Harm if the oath is major). Create an 8‑segment \textit{Covenant Bond} clock. When it fills, the covenant is considered witnessed by Mykkiel himself – breaking it after that point carries a Curse (GM chooses).
\end{itemize}

\subsection*{Advanced Rites (High Tier, 12+ XP)}

\begin{patronbox}{Divine Judgment [CURSE]}
\textbf{Cost:} 13 XP, Obligation 7 segments \\
\textbf{Effect:} Pronounce divine judgment on a target who has broken a covenant you have witnessed. The target suffers escalating penalties: first –1 die to all actions for the scene; if they persist, –2 dice; if they still defy, they take Harm 2 (Radiance) and are marked with a visible brand (e.g., a glowing handprint on their face). The judgment lasts until the target atones – publicly acknowledges their breach and makes restitution. \\
\textbf{Strategic note:} Use this on a villain who has betrayed a treaty, a spouse who has broken a marriage vow, or a merchant who has sold adulterated goods. The judgment is not secret – witnesses will know. Mykkiel’s justice is public.
\end{patronbox}

\begin{patronbox}{The Chosen Legion [WARD]}
\textbf{Cost:} 14 XP, Obligation 8 segments \\
\textbf{Effect:} Sanctify a large area (a fortress, a temple, a march border) as consecrated ground for the faithful. Within the zone, believers gain +1 die to defence rolls; unbelievers suffer –1 die. Supernatural entities who are not party to the covenant (e.g., demons, fae, outsiders) must test Spirit+Resolve (DV 4) each round or be unable to act. The zone lasts for one month or until the covenant is broken. \\
\textbf{Strategic note:} This is a campaign‑level rite. Use to protect a base of operations, to secure a holy site, or to create a safe haven during a war. The zone requires an anchor – a standing stone, a consecrated banner, a living covenant keeper. Destroy the anchor, and the zone collapses.
\end{patronbox}

\subsection*{GM Guidance – Mykkiel’s Intrusions}

When Mykkiel’s Obligation exceeds double capacity, choose or roll:
\begin{enumerate}
    \item A covenant you thought fulfilled is re‑examined by a third party – a rival arbiter, a fae court, an Inquisitor. They find a technical violation. You must either defend the covenant (social combat, DV 4) or pay a penance (1 Obligation or a minor sacrifice).
    \item Your signature appears on a document you have never seen – a contract, a treaty, a marriage licence. The handwriting is yours, but the memory is not. Someone has forged your bond.
    \item The scales of justice literally appear before you – in a dream, as a hallucination, as a physical object. They weigh two things: a good deed you have done and a bad one. If the bad outweighs the good, you lose 1 Boon (converted to a GM SB). If the good outweighs the bad, you gain +1 Boon. The scales are never balanced.
\end{enumerate}

\noindent \textbf{Runekeeper Note – The Codex of Seals:}
Mykkiel’s Codex is a collection of sealed documents, each containing a single Rite written in iron gall ink. To cast a Rite, you must break the seal – which consumes the document. This means Runekeepers of Mykkiel carry dozens of scrolls (or a heavy ledger of tear‑out pages). Losing the Codex means losing not just the Rites, but the *authority* to cast them – you must petition Mykkiel for new seals, which costs Obligation.

\clearpage

% ----------------------------------------------------------------------
% Raéyn – Mistress of the Sea (Tides & Change)
% ----------------------------------------------------------------------
\section{Raéyn – Mistress of the Sea}
\label{patron:raeyn}

\subsection*{Biography and Alignment}

\textit{“Mark the tide, name your course, and trust the wave‑road. But speak ill of Khemesh, and even I may let the deep take you.”}

Raéyn is the tempestuous goddess of the sea, the restless tide that carries news between shores and the promise of change between lives. She is mother to all who sail, her voice the wind that fills sails and her moods the storms that test every mariner’s resolve. But Raéyn’s heart is torn by her greatest tragedy: her son Khemesh, the Kraken of the Depths, who embodies the crushing inevitability of the ocean’s dark heart. Where Raéyn brings change and opportunity, Khemesh brings the final, inescapable pressure that grinds all things to nothing. Sailors pray to Raéyn for safe passage and favourable winds, but whisper Khemesh’s name when seeking to lay the dead to rest beneath the waves.

\textbf{Vignette – A Captain’s Bargain:}
\epigraph{“We were becalmed for nine days, the crew eyeing each other with thirst‑madness. I knelt on the bow and poured my last bottle of fine Kahfagian wine into the sea – not as a bribe, but as a toast. ‘To the storms that test us,’ I said. ‘To the shoals that teach us.’ An hour later, a squall came from nowhere – not a killer, just a push. It filled our sails and drove us to port by midnight. I’ve since learned: Raéyn does not want slaves. She wants sailors who respect the sea for what it is – a lover, a teacher, and a grave.”}{\textit{— Captain Sera of the Wave‑Breaker}}

\textbf{What They Represent:} Tides, change, sea travel, storms, the dual nature of the ocean (giver and taker). Raéyn is not a goddess of calm waters – she is a goddess of *passage*. She tests those who sail her, rewards those who respect her, and mourns those who drown. Her followers understand that every harbour is temporary, every voyage a risk, and every return a blessing.

\textbf{Cultural Variants:}
\begin{itemize}
    \item \textbf{Kahfagia:} \textit{the Lantern‑Tide}. Pilots light a beacon to Raéyn before every night voyage. Her symbol is a wave curling into a spiral.
    \item \textbf{Linn:} \textit{the Whale‑Mother}. Raéyn is the great she‑whale who carries the drowned to the afterlife in her belly. Funerals at sea involve a song and a single white stone dropped overboard.
    \item \textbf{Zakov:} \textit{the Salt‑Queen}. Pirates carve her face on their figureheads – one eye open (fair seas), one eye closed (storms ahead). She is honoured and feared.
    \item \textbf{The Inquisitor Prime’s View:} Raéyn is a wild, untamed power, not a heretic. The Inquisition tolerates her worship among coastal folk but watches for cults that “drown unbelievers” – those are Malachai’s work, not Raéyn’s.
\end{itemize}

\subsection*{Playing a Follower of Raéyn}

\textbf{Strategy:} You are the party’s navigator, weather‑reader, and master of maritime environments. Raéyn’s rites excel at travel on or near water, weather manipulation, and the ebb and flow of tides. In combat, you control movement and environment – creating waves that knock enemies prone, summoning fog to hide an escape, or calling a favourable wind for a fleeing ship. Out of combat, you are the one who knows the currents, who can predict storms, who can read a sailor’s heart.

\textbf{Suggested Acts of Obligation (to clear debt):}
\begin{itemize}
    \item Pour a libation into the sea (wine, beer, fresh water) before a voyage, asking for nothing in return.
    \item Rescue a drowning creature – a sailor, a fisher, even a seabird – without expectation of reward.
    \item Return a piece of flotsam to the water. It may have washed ashore; the sea wants it back.
    \item Walk the shore at dawn or dusk and listen. If you hear the waves speaking a name, speak it back. Raéyn will remember.
\end{itemize}

\subsection*{Rites – Strategic Use}

\textbf{Low Rites}
\begin{itemize}
    \item \textit{Tidemark’s Blessing} (4 XP): Treat slick, swaying, or water‑slicked footing as stable for you this scene. Gain +1 die on boarding, balance, or shipboard movement. Use when fighting on a wet deck or crossing a slippery reef.
    \item \textit{Whispering Currents} (5 XP): Learn the safest near‑term route across water or coastline (reefs, eddies, patrols) or gain +1 die to navigation checks for the scene. Use to avoid a naval ambush or to find a hidden smuggler’s cove.
\end{itemize}

\textbf{Standard Rites}
\begin{itemize}
    \item \textit{The Changing Tide} (8 XP): Bias currents and water levels in a zone. Those moving with the tide gain +1 die; those moving against suffer –1 die. Small craft must test to hold position. Use to favour your ship in a chase or to strand an enemy vessel on a sandbar.
    \item \textit{Wave‑Road Blessing} (7 XP): Consecrate a wave‑road between two visible points. Allies gain +2 dice to travel, evade, or carry actions at sea; designated pursuers suffer –1 die to intercept. Use for a critical voyage – a courier mission, a smuggling run, a desperate escape.
\end{itemize}

\subsection*{Advanced Rites (High Tier, 12+ XP)}

\begin{patronbox}{The Storm‑Queen’s Hand [AREA]}
\textbf{Cost:} 13 XP, Obligation 7 segments \\
\textbf{Effect:} Shape a storm‑band over a zone (sea or shore). At casting, choose two modes; you may switch one mode once per scene:
\begin{itemize}
    \item \textbf{Propulsion:} Allied vessels gain +1 band of movement per beat (or +1 Effect to manoeuvres).
    \item \textbf{Concealment:} Veil of rain and spray; ranged targeting impaired; enemies suffer –1 die to sight‑based perception.
    \item \textbf{Smite:} Once per beat, lash a target with wave or lightning as an [AREA] hazard (DV 3, Harm 1, Push to upgrade).
\end{itemize}
The storm lasts for the scene. When it ends, the caster marks 1 Fatigue (eyes full of salt). \\
\textbf{Strategic note:} Use this to turn a naval battle, to cover an amphibious landing, or to escape a blockade. Raéyn gives you a taste of her power – but the storm answers to her, not you.
\end{patronbox}

\begin{patronbox}{The Mother’s Wrath [CURSE]}
\textbf{Cost:} 14 XP, Obligation 7 segments \\
\textbf{Effect:} Curse a target who has wronged you – a pirate who burned your ship, a merchant who sold you a leaky hull, a noble who polluted a sacred bay. For one session, the target suffers –2 dice to all maritime and weather‑related rolls. At sea, they must test Spirit+Resolve (DV 4) each day or suffer Harm 1 (Weather – a sudden squall, a rogue wave). The curse can be lifted if the target makes a sincere apology to Raéyn (a libation, a gift, a public admission of fault). \\
\textbf{Strategic note:} This is divine vengeance, not petty spite. Use it only against those who have truly violated the sea’s law – polluters, betrayers of sailors, those who drown the innocent. Raéyn does not assist in personal feuds.
\end{patronbox}

\subsection*{GM Guidance – Raéyn’s Intrusions}

When Raéyn’s Obligation exceeds double capacity, choose or roll:
\begin{enumerate}
    \item A sudden fog rolls in – not hostile, but disorienting. All navigation rolls are at –2 dice for the scene. The fog carries whispers of drowned sailors. They do not attack; they merely… advise.
    \item The tide turns early, stranding a vessel you are on (or a vessel you need). You must either wait six hours (advance a Supply clock) or find a creative way to refloat (Body+Strength, DV 4).
    \item A child of Raéyn – a dolphin, a seal, a sea eagle – follows the party for a day. It is a messenger. It will deliver one short message to the nearest follower of Raéyn, but only if you feed it something shiny.
\end{enumerate}

\noindent \textbf{Cantor Variant – The Shanty of the Deep:}
Cantors of Raéyn sing work songs, sea shanties, wordless melodies that mimic wind and wave. Their Performance Maneuver does not shift ally Position; instead, it allows the party to ignore one level of environmental penalty (storm, fog, choppy water) for one exchange. The Cantor must stamp their foot or clap in rhythm – the sound anchors the effect.

\clearpage

\clearpage

\chapter{Rituals and Components – When You Lack Focus}

A \textbf{ritual} is any magical working that takes longer than one action. Normally, Runekeepers use their \textit{Codex}, Invokers use a \textit{Symbol}, and Cantors use their \textit{voice} (no focus needed). When you lack the usual focus, you must supply \textbf{components}.

\section{The Core Substitution Rule}

\begin{ritualnote}
\textbf{If you lack your focus:}
\begin{itemize}
    \item \textbf{Runekeeper (no Codex):} You must have a written record of the Rite – even a scrap of paper. Copying from memory requires a \texttt{Wits+Lore} test (DV 4); failure means the ritual fizzles and you mark 1 Fatigue.
    \item \textbf{Invoker (no Symbol):} You need a physical object tied to the Patron (see component examples below). Using it consumes the object.
    \item \textbf{Cantor (no audience):} You cannot use Performance Maneuvers. Casting Rites increases DV by +1. An enchanted instrument can substitute, but its curse still applies.
\end{itemize}
\textbf{Time increase:} +1 step (e.g., Some Time → Significant Time). \\
\textbf{DV increase:} +1.
\end{ritualnote}

\section{Component Examples by Patron Theme}

\begin{center}
\begin{longtable}{p{3cm} p{4cm} p{4cm}}
\caption{Components for Missing Focus}\\
\toprule
Patron (theme) & Component (consumed) & Where to find \\
\midrule
Aveh (Freedom) & Broken chain link, bird feather & Anywhere oppression exists \\
Khemesh (Depth) & Seashell from deep trench, drowned coin & Coastlines, shipwrecks \\
Palinode (Performance) & Unseen sheet music, a tear & Theatres, crowded squares \\
Morag (Bargains) & Signed contract, faerie hair & Markets, crossroads \\
Inquisitor Prime (Purity) & Drop of silver, sinner’s confession & Churches, prisons \\
Traveler (Ways) & Worn bootlace, pinch of road dust & Crossroads, mile markers \\
\bottomrule
\end{longtable}
\end{center}

\section{Ritual Without Focus – Step by Step}
\label{sec:ritual-steps}
\begin{enumerate}
    \item State your intent.
    \item Declare your missing focus (Codex, Symbol, or audience).
    \item Narrate gathering a component (use the table above or invent one).
    \item Time increases by one step.
    \item DV increases by +1.
    \item Roll (usually \texttt{Wits+Arcana} or \texttt{Spirit+Lore}).
    \item \textbf{On a miss:} Component destroyed, GM gains +1 SB. \\
          \textbf{On a partial:} Ritual works, component consumed, and you mark 1 Fatigue.
\end{enumerate}

\subsection{Example – Invoker Without Symbol}
Kestra (Invoker of the Pale Shepherd) wants to cast \textit{Rite of the Hidden Path} to escape a collapsing cave, but she left her Symbol (a grave‑stone chip) in her pack.

\begin{itemize}
    \item \textbf{GM:} “You need a component – something tied to passage or thresholds.”
    \item \textbf{Kestra:} “I have a lock of hair from a friend who died. Does that work?”
    \item \textbf{GM:} “Yes. Ritual takes 5 minutes instead of 1. DV is 4 (base 3+1).”
    \item \textbf{Roll:} Success! The hair burns away, and a hidden fissure opens.
\end{itemize}

\clearpage

\chapter{Quick Reference Tables}

\section{Ritual Substitution Summary}
\begin{longtable}{p{2.5cm} p{2.5cm} p{2.5cm} p{2.5cm} p{3cm}}
\toprule
Path & Normal Focus & Missing → Component & Time Increase & DV Increase \\
\midrule
Runekeeper & Codex & Written scrap (copied from memory) & +1 step & +1 \\
Invoker & Symbol & Patron‑themed object & +1 step & +1 \\
Cantor & (voice – no focus) & Enchanted instrument or audience & +1 step & +0 (audience willing) / +1 (hostile) \\
\bottomrule
\end{longtable}

\section{Obligation at a Glance}
\begin{itemize}
    \item Capacity = Spirit + Presence.
    \item Over capacity → mark Fatigue per segment over.
    \item Double capacity → Patron Intrusion (GM chooses from patron box).
    \item Clear by: Act of Service, Ritual Purification, or GM Fiat.
\end{itemize}

\clearpage

% ----------------------------------------------------------------------
% Raéyn – Mistress of the Sea (Tides & Change)
% ----------------------------------------------------------------------
\section{Raéyn – Mistress of the Sea}
\label{patron:raeyn}

\subsection*{Biography and Alignment}

\textit{“Mark the tide, name your course, and trust the wave‑road. But speak ill of Khemesh, and even I may let the deep take you.”}

Raéyn is the tempestuous goddess of the sea, the restless tide that carries news between shores and the promise of change between lives. She is mother to all who sail, her voice the wind that fills sails and her moods the storms that test every mariner’s resolve. But Raéyn’s heart is torn by her greatest tragedy: her son Khemesh, the Kraken of the Depths, who embodies the crushing inevitability of the ocean’s dark heart. Where Raéyn brings change and opportunity, Khemesh brings the final, inescapable pressure that grinds all things to nothing. Sailors pray to Raéyn for safe passage and favourable winds, but whisper Khemesh’s name when seeking to lay the dead to rest beneath the waves.

\textbf{Vignette – A Captain’s Bargain:}
\epigraph{“We were becalmed for nine days, the crew eyeing each other with thirst‑madness. I knelt on the bow and poured my last bottle of fine Kahfagian wine into the sea – not as a bribe, but as a toast. ‘To the storms that test us,’ I said. ‘To the shoals that teach us.’ An hour later, a squall came from nowhere – not a killer, just a push. It filled our sails and drove us to port by midnight. I’ve since learned: Raéyn does not want slaves. She wants sailors who respect the sea for what it is – a lover, a teacher, and a grave.”}{\textit{— Captain Sera of the Wave‑Breaker}}

\textbf{What They Represent:} Tides, change, sea travel, storms, the dual nature of the ocean (giver and taker). Raéyn is not a goddess of calm waters – she is a goddess of *passage*. She tests those who sail her, rewards those who respect her, and mourns those who drown. Her followers understand that every harbour is temporary, every voyage a risk, and every return a blessing.

\textbf{Cultural Variants:}
\begin{itemize}
    \item \textbf{Kahfagia:} \textit{the Lantern‑Tide}. Pilots light a beacon to Raéyn before every night voyage. Her symbol is a wave curling into a spiral.
    \item \textbf{Linn:} \textit{the Whale‑Mother}. Raéyn is the great she‑whale who carries the drowned to the afterlife in her belly. Funerals at sea involve a song and a single white stone dropped overboard.
    \item \textbf{Zakov:} \textit{the Salt‑Queen}. Pirates carve her face on their figureheads – one eye open (fair seas), one eye closed (storms ahead). She is honoured and feared.
    \item \textbf{The Inquisitor Prime’s View:} Raéyn is a wild, untamed power, not a heretic. The Inquisition tolerates her worship among coastal folk but watches for cults that “drown unbelievers” – those are Malachai’s work, not Raéyn’s.
\end{itemize}

\subsection*{Playing a Follower of Raéyn}

\textbf{Strategy:} You are the party’s navigator, weather‑reader, and master of maritime environments. Raéyn’s rites excel at travel on or near water, weather manipulation, and the ebb and flow of tides. In combat, you control movement and environment – creating waves that knock enemies prone, summoning fog to hide an escape, or calling a favourable wind for a fleeing ship. Out of combat, you are the one who knows the currents, who can predict storms, who can read a sailor’s heart.

\textbf{Suggested Acts of Obligation (to clear debt):}
\begin{itemize}
    \item Pour a libation into the sea (wine, beer, fresh water) before a voyage, asking for nothing in return.
    \item Rescue a drowning creature – a sailor, a fisher, even a seabird – without expectation of reward.
    \item Return a piece of flotsam to the water. It may have washed ashore; the sea wants it back.
    \item Walk the shore at dawn or dusk and listen. If you hear the waves speaking a name, speak it back. Raéyn will remember.
\end{itemize}

\subsection*{Rites – Strategic Use}

\textbf{Low Rites}
\begin{itemize}
    \item \textit{Tidemark’s Blessing} (4 XP): Treat slick, swaying, or water‑slicked footing as stable for you this scene. Gain +1 die on boarding, balance, or shipboard movement. Use when fighting on a wet deck or crossing a slippery reef.
    \item \textit{Whispering Currents} (5 XP): Learn the safest near‑term route across water or coastline (reefs, eddies, patrols) or gain +1 die to navigation checks for the scene. Use to avoid a naval ambush or to find a hidden smuggler’s cove.
\end{itemize}

\textbf{Standard Rites}
\begin{itemize}
    \item \textit{The Changing Tide} (8 XP): Bias currents and water levels in a zone. Those moving with the tide gain +1 die; those moving against suffer –1 die. Small craft must test to hold position. Use to favour your ship in a chase or to strand an enemy vessel on a sandbar.
    \item \textit{Wave‑Road Blessing} (7 XP): Consecrate a wave‑road between two visible points. Allies gain +2 dice to travel, evade, or carry actions at sea; designated pursuers suffer –1 die to intercept. Use for a critical voyage – a courier mission, a smuggling run, a desperate escape.
\end{itemize}

\subsection*{Advanced Rites (High Tier, 12+ XP)}

\begin{patronbox}{The Storm‑Queen’s Hand [AREA]}
\textbf{Cost:} 13 XP, Obligation 7 segments \\
\textbf{Effect:} Shape a storm‑band over a zone (sea or shore). At casting, choose two modes; you may switch one mode once per scene:
\begin{itemize}
    \item \textbf{Propulsion:} Allied vessels gain +1 band of movement per beat (or +1 Effect to manoeuvres).
    \item \textbf{Concealment:} Veil of rain and spray; ranged targeting impaired; enemies suffer –1 die to sight‑based perception.
    \item \textbf{Smite:} Once per beat, lash a target with wave or lightning as an [AREA] hazard (DV 3, Harm 1, Push to upgrade).
\end{itemize}
The storm lasts for the scene. When it ends, the caster marks 1 Fatigue (eyes full of salt). \\
\textbf{Strategic note:} Use this to turn a naval battle, to cover an amphibious landing, or to escape a blockade. Raéyn gives you a taste of her power – but the storm answers to her, not you.
\end{patronbox}

\begin{patronbox}{The Mother’s Wrath [CURSE]}
\textbf{Cost:} 14 XP, Obligation 7 segments \\
\textbf{Effect:} Curse a target who has wronged you – a pirate who burned your ship, a merchant who sold you a leaky hull, a noble who polluted a sacred bay. For one session, the target suffers –2 dice to all maritime and weather‑related rolls. At sea, they must test Spirit+Resolve (DV 4) each day or suffer Harm 1 (Weather – a sudden squall, a rogue wave). The curse can be lifted if the target makes a sincere apology to Raéyn (a libation, a gift, a public admission of fault). \\
\textbf{Strategic note:} This is divine vengeance, not petty spite. Use it only against those who have truly violated the sea’s law – polluters, betrayers of sailors, those who drown the innocent. Raéyn does not assist in personal feuds.
\end{patronbox}

\subsection*{GM Guidance – Raéyn’s Intrusions}

When Raéyn’s Obligation exceeds double capacity, choose or roll:
\begin{enumerate}
    \item A sudden fog rolls in – not hostile, but disorienting. All navigation rolls are at –2 dice for the scene. The fog carries whispers of drowned sailors. They do not attack; they merely… advise.
    \item The tide turns early, stranding a vessel you are on (or a vessel you need). You must either wait six hours (advance a Supply clock) or find a creative way to refloat (Body+Strength, DV 4).
    \item A child of Raéyn – a dolphin, a seal, a sea eagle – follows the party for a day. It is a messenger. It will deliver one short message to the nearest follower of Raéyn, but only if you feed it something shiny.
\end{enumerate}

\noindent \textbf{Cantor Variant – The Shanty of the Deep:}
Cantors of Raéyn sing work songs, sea shanties, wordless melodies that mimic wind and wave. Their Performance Maneuver does not shift ally Position; instead, it allows the party to ignore one level of environmental penalty (storm, fog, choppy water) for one exchange. The Cantor must stamp their foot or clap in rhythm – the sound anchors the effect.

\clearpage

% ----------------------------------------------------------------------
% General Rites – Exorcism, Consecration, and Other Sacred Works
% ----------------------------------------------------------------------
\chapter{General Rites – Works for Any Hand}

\epigraph{“Not every prayer needs a patron. Not every door needs a key. Some rites are older than the covenants – older than the gods themselves. These are the works of bare stone and steady breath.”}{\textit{— The Gray Wanderer}}

The rites in this chapter are not tied to any single Patron. They are the common magic of threshold, seal, and vigil – practices known to hedge‑priests, night‑watchmen, and village elders. They require no Symbol, no Codex, no Thiasos. What they require is \textbf{time}, \textbf{components}, and a \textbf{clock} – a visible measure of progress that anyone at the table can track.

\section{The Rite of Exorcism – Casting Out the Unwanted}

\textit{“Speak its name. Know its nature. Then tell it to leave – and mean it.”}

Exorcism is a contested ritual. You are not merely casting a spell; you are \textit{negotiating} with a possessing entity, and it does not want to go.

\begin{tcolorbox}[colback=blue!5, colframe=blue!70]
\textbf{Rite of Exorcism}
\begin{itemize}
    \item \textbf{Time:} Significant Time (phased, see below)
    \item \textbf{Components:} Salt (or iron filings), a vessel for the entity to depart into (a bowl of water, a candle flame, a bell), and the subject’s true name (or a close approximation)
    \item \textbf{Test:} Contested clock (size = maximum of 4 and 2 + entity’s Tier)
\end{itemize}
\end{tcolorbox}

\subsection*{The Exorcism Clock (6 or 8 segments)}

Both the exorcist and the entity work against each other. The clock represents the exorcist’s progress toward expulsion. On each of the exorcist’s turns, they choose one of three approaches and roll against the entity’s DV (usually the entity’s Tier + 2). The GM rolls for the entity to resist; if the entity meets or exceeds the exorcist’s successes, it ticks the \textbf{Possession Deepens} clock instead.

\subsection*{Approaches}

\begin{itemize}
    \item \textbf{Command (Presence + Command):} “I speak with authority. Leave this body.” DV = 3 + entity’s Tier. On success, tick Exorcism clock +1. On a miss, the entity ticks Possession Deepens and may harm the subject (Fatigue 1).
    \item \textbf{Lore (Wits + Lore):} “I name your origin and your weakness.” DV = 4 + entity’s Tier. On success, tick Exorcism clock +2 (you found a leverage point). On a miss, the entity gains +1 die on its next resistance.
    \item \textbf{Sacrifice (Spirit + Resolve):} “I offer something of myself in exchange for your departure.” DV = 3 + entity’s Tier. On success, tick Exorcism clock +1 and the exorcist marks 1 Fatigue (the cost is real). On a miss, the entity takes the sacrifice without leaving – the exorcist marks 1 Fatigue and the clock does not advance.
\end{itemize}

\subsection*{Entity’s Reactions}

At the end of each round (after the exorcist’s action and any allied help), the GM may have the entity attempt a **Response** (spend 1 SB to escalate). Choose:
\begin{itemize}
    \item \textbf{Whispers:} The subject speaks in a voice not their own, revealing a secret about one of the exorcists (GM gains 1 SB).
    \item \textbf{Manifestation:} The entity briefly appears – ghostly, insectile, shadowy – and attacks a nearby object or bystander (tick a relevant clock).
    \item \textbf{Assault:} The subject makes a violent lunge at the exorcist (Body + Melee, DV 3, Harm 1 if it hits).
\end{itemize}

\subsection*{Resolution}

\begin{itemize}
    \item \textbf{Exorcism clock fills first:} The entity is expelled. It may leave into the prepared vessel, which then cracks or shatters (the rite is complete). The subject is exhausted (mark Fatigue 2) but free.
    \item \textbf{Possession Deepens clock fills first:} The entity takes deeper root. The subject suffers Harm 2 (Spiritual) and gains a permanent tell (e.g., a second shadow, a cold touch). The exorcism fails; the entity may now act through the subject at will once per session until a more powerful rite is attempted.
\end{itemize}

\subsection*{Help and Teamwork}

Up to three assistants can aid the exorcist. Each chooses one of the three approaches and rolls separately; their successes add to the exorcist’s total for that round. However, each assistant who rolls a 1 generates an SB that the GM may spend on the entity’s Reactions.

\clearpage

\section{The Rite of Consecration – Making Ground Sacred}

\textit{“This place is now a boundary. What passes must be willing. What intrudes must be invited – or destroyed.”}

Consecration turns a space into a temporary sanctuary. It can be used to protect a campsite, to prepare a ritual circle, or to deny entry to hostile supernatural beings.

\begin{tcolorbox}[colback=blue!5, colframe=blue!70]
\textbf{Rite of Consecration}
\begin{itemize}
    \item \textbf{Time:} Some Time (10–30 minutes, depending on area)
    \item \textbf{Components:} Salt, ash, or chalk; a blessed object (a relic, a symbol, a bell); and a spoken declaration of intent
    \item \textbf{Clock:} \textit{Sanctity} (6 segments) – the strength of the ward
\end{itemize}
\end{tcolorbox}

\subsection*{Procedure}

The caster walks the perimeter of the area, marking it with the chosen component (salt line, chalk sigils, ash smudges). At each cardinal point (or significant threshold), they make a \textbf{Spirit + Lore} test against DV 3.

\begin{itemize}
    \item \textbf{Success:} Tick Sanctity clock +1.
    \item \textbf{Partial:} Tick Sanctity +1, but the caster marks 1 Fatigue.
    \item \textbf{Miss:} No tick; the component used at that point is wasted (GM may force a choice: continue with insufficient markings, or abandon the rite).
\end{itemize}

When the Sanctity clock reaches 4 segments, the area gains the \textbf{[WARD]} tag against \textit{one type} of intruder (e.g., undead, fae, spirits, animals). At 6 segments, it gains \textbf{[WARD]} against all supernatural entities not keyed to the consecration (the caster and their allies can pass freely; enemies must test Spirit+Resolve (DV 4) or be blocked).

\subsection*{Maintaining Consecration}

The consecration lasts for one night or until the perimeter is broken (a line scuffed, a sigil smeared, a threshold crossed without permission). The caster can renew it by spending another Some Time and repeating the rite – but the Sanctity clock resets to 0 unless they spend 2 XP to anchor it permanently (rare).

\subsection*{Consecration Gimmick – Lingering Echo}

When the consecration ends (by time or by breach), the caster may roll 1d6. On a 4+, the residue of the rite remains for one more scene – not a full ward, but a **Boon** to anyone who stays inside: +1 die to resist fear or supernatural influence.

\clearpage

\section{The Rite of Sealing – Locking a Threshold}

\textit{“This door does not open. This path does not lead. This gate remembers its purpose.”}

Sealing is a simpler, faster version of the Sealed Gate’s rites – open to anyone with a bit of salt and a strong will.

\begin{tcolorbox}[colback=blue!5, colframe=blue!70]
\textbf{Rite of Sealing}
\begin{itemize}
    \item \textbf{Time:} One minute (one action in combat; several minutes outside)
    \item \textbf{Components:} A mark (chalk, blood, charcoal) or a physical barrier (a tied rope, a nailed board)
    \item \textbf{Clock:} \textit{Seal Integrity} (4 segments) – how many attempts to break it
\end{itemize}
\end{tcolorbox}

\subsection*{Procedure}

The caster declares a threshold (a door, a window, a path between two stones) and marks it. Roll \textbf{Spirit + Resolve} (DV 3). On success, the seal is set with 3 Integrity segments. On partial, 2 segments. On miss, 1 segment – the GM may also note that the seal is imperfect (a faint glow, a smell of ozone).

\subsection*{Breaking the Seal}

Any creature attempting to cross the sealed threshold must first \textbf{test Body + Resolve} (DV = remaining Integrity segments). On success, they cross and reduce Seal Integrity by 1. On failure, they cannot cross this round and may be knocked back (GM’s choice). Physical force (a battering ram, a spell of shattering) can also reduce Integrity – treat as an attack against the seal (DV 4, each success reduces Integrity by 1).

\subsection*{Sealing Gimmick – False Seal}

If the caster rolls a critical success (at least 1 die showing 10 and total successes ≥ DV by 2+), the seal appears stronger than it is. The first creature to test against it believes the remaining Integrity is 2 higher than actual (the GM secretly notes the real value). After that first attempt, the illusion shatters.

\clearpage

\section{The Rite of Vigil – Watching Through the Night}

\textit{“Sleep is a small death. I will stand guard between the living and the dark.”}

Vigil is not a ritual of power but a ritual of endurance. It is used to protect a sleeping group, to watch for ambush, or to hold a line against subtle threats (nightmares, spirits, vermin).

\begin{tcolorbox}[colback=blue!5, colframe=blue!70]
\textbf{Rite of Vigil}
\begin{itemize}
    \item \textbf{Time:} An entire night (8 hours of watch)
    \item \textbf{Components:} A flame (candle, lantern, torch) and a weapon or symbol held in the lap
    \item \textbf{Clock:} \textit{Vigil Fatigue} (4 + caster’s Body) – each segment is an hour of wakefulness
\end{itemize}
\end{tcolorbox}

\subsection*{Procedure}

At the start of the night, the caster rolls \textbf{Body + Endurance} (DV 3). On success, their Vigil Fatigue clock starts at 0. On partial, they begin with 1 Fatigue already marked (they were tired). On miss, they start with 2 Fatigue and the GM may introduce a minor nocturnal complication (an owl, a creaking branch).

Every hour (or every scene interruption), the GM may call for a new Body + Endurance test (DV increases by 1 each time, to a maximum of 6). On a success, the vigil holds. On a partial, the caster marks 1 Fatigue and the GM may advance an \textit{Approaching Threat} clock by 1. On a miss, the caster falls asleep (or suffers a vision) – the vigil fails, and the party is vulnerable.

\subsection*{Vigil Gimmick – The Watcher’s Boon}

If the caster completes the full night without falling asleep (Vigil Fatigue never filled), all allies who rested gain +1 die to their first action of the next day (a boon of alertness). The caster also clears 2 Fatigue from the vigil (the satisfaction of duty well done).

\clearpage

\section{The Rite of the Weary Road – Purifying Exhaustion}

\textit{“You have walked too far. You have carried too much. Sit. Breathe. Let the road fall from your shoulders.”}

A folk rite of travellers, used to shake off the spiritual weight of a long journey. It does not heal wounds but eases the mind and the spirit.

\begin{tcolorbox}[colback=blue!5, colframe=blue!70]
\textbf{Rite of the Weary Road}
\begin{itemize}
    \item \textbf{Time:} One hour (cannot be done in a rush)
    \item \textbf{Components:} Clean water, a pinch of salt, and a few words spoken to the nearest landmark (a tree, a stone, a signpost)
    \item \textbf{Clock:} \textit{Spiritual Weight} (equal to the number of days travelled since last rest)
\end{itemize}
\end{tcolorbox}

\subsection*{Procedure}

The caster leads the group in a short meditation. Each participant rolls \textbf{Spirit + Survival} (DV = current Spiritual Weight). On success, they clear 1 Fatigue (mental only, not physical) and reduce the Spiritual Weight clock by 1 for the next person. On partial, they clear 1 Fatigue but the weight does not reduce. On miss, they gain 1 Fatigue (the meditation brought up dark memories).

After each participant has rolled, the rite ends. Anyone who succeeded at least once may also re‑roll a failed social or investigation roll from the previous day (the clarity of the rite gave them a new perspective).

\subsection*{Weary Road Gimmick – Shared Burden}

If the group rolls cooperatively (each assisting another), the total number of successes across all participants reduces the Spiritual Weight clock for everyone by that amount, and each participant clears 1 Fatigue regardless of individual rolls. However, if any participant rolls a 1, the GM gains 1 SB to spend on a flashback complication – something from the road they thought they had left behind.

\clearpage

\section{Quick Reference – General Rites}

\begin{longtable}{p{3cm} p{2.5cm} p{2.5cm} p{5cm}}
\toprule
Rite & Time & Clock & Gimmick \\
\midrule
Exorcism & Significant (phased) & Contested (6–8) & Entity can resist; Possession Deepens clock \\
Consecration & Some Time & Sanctity (6) & Lingering Echo (Boon after end) \\
Sealing & 1 minute & Seal Integrity (4) & False Seal (illusory strength) \\
Vigil & 8 hours & Vigil Fatigue (4+Body) & Watcher’s Boon (+1 die next day) \\
Weary Road & 1 hour & Spiritual Weight (days travelled) & Shared Burden (cooperative reduction) \\
\bottomrule
\end{longtable}

\clearpage

% ----------------------------------------------------------------------
% The Cantor’s Cult – How Song Feeds the Divine
% ----------------------------------------------------------------------
\chapter{The Cantor’s Cult – The Subtle Cycle}
\epigraph{“A Runekeeper writes a Rite. An Invoker borrows a Symbol. But a Cantor sings, and a crowd listens – and in that listening, something wakes.”}{\textit{— The Gray Wanderer}}

Cantors are the most public of the three paths. They cannot hide their magic; a hymn is not a whisper. But publicity is not a weakness – it is a \textit{fuel source}. Every audience, every captivated crowd, every tear shed to a melody becomes a kind of offering. The Patrons notice. And from that noticing, a cycle begins.

\section{The Unwritten Chorus – How Cantors Feed Patrons}

A Runekeeper pays Obligation through service, sacrifice, or ritual. An Invoker pays through Symbol maintenance and the risk of \textit{Crack the Seal}. A Cantor pays through \textit{attention}.

When a Cantor performs – especially a Resonant Rite or a Performance Maneuver before a crowd of three or more – they generate \textbf{Resonance}. The GM tracks Resonance as an invisible resource (a 4‑segment clock, not shown to players). Each time a Cantor rolls a 1 on a Performance roll in front of an audience, tick Resonance +1. Each time a crowd reacts emotionally (gasps, cheers, weeping), tick Resonance +1.

\subsection*{Resonance Payoff}

When the Resonance clock fills (4 segments), the Cantor’s Patron (or the nearest listening Patron, if the Cantor is unsworn) gains a \textbf{Debt of Attention}. The GM may convert this Debt into:

\begin{itemize}
    \item A free \textbf{Act of Service} for the Cantor (clear 1 Obligation without roleplay).
    \item A \textbf{Dream} sent to a Runekeeper of the same Patron somewhere in the world (the dream contains a fragment of the crowd’s emotion – useful for a Rite component).
    \item A \textbf{Waking Echo} – a brief manifestation of the Patron’s presence (a candle lighting itself, a bell ringing unpulled) that can be used as a Symbol component for an Invoker of that Patron who is nearby.
\end{itemize}

The cycle is subtle: the Cantor’s audience feeds the Patron, the Patron rewards a Runekeeper (who owes Obligation), and the Runekeeper may teach a Rite to the Cantor in thanks. No one sees the whole loop; each only sees their part.

\section{Recruitment – How Cantors Attract Runekeepers}

A Runekeeper is an archivist, a scholar, a keeper of written Rites. They rarely proselytise. Cantors, however, attract followers – a singer who moves a crowd may also move a young scholar to ask, “How does that work? Can I write it down?”

\subsection*{The Recruitment Clock (6 segments)}

When a Cantor performs a significant public rite (a funeral hymn, a battle chant, a festival song) before an audience that includes a potential Runekeeper – a scribe, a temple acolyte, a curious merchant – the GM may start a \textbf{Recruitment} clock.

\begin{itemize}
    \item \textbf{Tick +1:} The audience member asks a question about the magic.
    \item \textbf{Tick +2:} They request a private meeting or a written copy of a verse.
    \item \textbf{Tick +3:} They offer to transcribe the Cantor’s songs into a Codex (the first step toward becoming a Runekeeper).
    \item \textbf{Fill:} The audience member formally asks to study under the Cantor – not to learn to sing, but to learn the \textit{mechanics} of the Patron’s power. They become a Follower (Cap 1) with a single Rite; within a few sessions, they may become a full Runekeeper PC or NPC.
\end{itemize}

The Cantor does not control this clock directly – it advances when the GM feels the story warrants it. But the Cantor can accelerate it by \textbf{spending 1 Boon} to create an opportunity for a meaningful interaction (e.g., “I catch the scribe’s eye during the chorus and hold it for a breath longer than necessary.”)

\subsection*{The Reverse – Runekeepers as Cantor Patrons}

Occasionally, a Runekeeper’s Codex contains a Song that no Cantor has ever sung. If the Runekeeper teaches it to a Cantor, the Runekeeper gains 1 Obligation clearance (as if they had performed an Act of Service) and the Cantor gains a new Rite. The Runekeeper’s Patron may also send a dream of gratitude – not to the Runekeeper, but to the Cantor, reinforcing the cycle.

\section{The Cult as Middle Ground – Cantors, Cults, and the Inquisition}

A Cantor’s following can become a cult – not necessarily sinister, but always intense. Cults of Palinode are performance troupes that never disband; cults of Malachai are secret societies of cursed individuals who share their burdens; cults of the Witness are investigative circles that share uncovered truths.

\subsection*{Cult Mechanics (Optional)}

If a Cantor gains three or more dedicated followers (Cap 1 each), the group may be treated as a \textbf{Cult Asset} (Standard, 8 XP value). Benefits include:

\begin{itemize}
    \item \textbf{Free Resonance:} Once per session, the Cantor may declare a crowd scene and automatically tick Resonance +1 (no roll needed).
    \item \textbf{Rumour Network:} The cult spreads news; the Cantor gains +1 die to Streetwise or Investigation rolls in any city where the cult has a chapter.
    \item \textbf{Shelter:} The cult can hide the party from authorities for up to a week, but each day of hiding ticks a \textit{Heat} clock (4 segments). When Heat fills, the Inquisition (or local authorities) raids the safehouse.
\end{itemize}

\subsection*{The Inquisition’s View}

Inquisitors of the Prime and witch‑hunters of other faiths watch Cantors carefully. A Cantor with a large following is a potential heretic – not because they are evil, but because crowds are unpredictable. A Cantor can reduce Inquisition attention by:

\begin{itemize}
    \item Performing for a temple (instead of a tavern) – gain a \textbf{Sanctioned Performer} token (reduces Heat gain by 1 for a scene).
    \item Teaching a Rite to an Inquisitor’s agent (the agent becomes a mole – the Cantor gains 2 Boons, but the mole reports everything).
    \item Publicly disowning a follower who becomes too zealous (clears 1 Obligation but loses that follower permanently).
\end{itemize}

\section{Example – A Cycle in Play}

\textit{Kestra (Cantor of Palinode) performs at a Silkstrand festival. She rolls a 1 on her Performance Maneuver – the GM ticks Resonance. The crowd laughs at a wrong note. A scribe in the front row (a potential Runekeeper) writes something down. Later, she sings a Resonant Rite – the Resonance clock fills.}

\textit{The GM decides that Palinode sends a Dream Echo to a Runekeeper named Thera in a distant city. Thera dreams of a silk loom weaving itself. Inspiration strikes – Thera writes a new Rite of Binding Threads. Thera owes 1 Obligation to Palinode for the dream.}

\textit{Months later, Thera meets Kestra by chance. They compare notes. Thera teaches Kestra the new Rite. Kestra now has a Rite she never would have learned alone. The cycle completes.}

\section{GM Guidance – Pacing the Cycle}

\begin{itemize}
    \item \textbf{Resonance clocks should fill once every 2–3 sessions.} If they fill faster, the Patrons become too present (weakening the mystery). If slower, the Cantor feels disconnected.
    \item \textbf{Recruitment clocks are optional.} Not every Cantor wants followers. A player may refuse the clock by having their Cantor deliberately avoid mentoring.
    \item \textbf{The Inquisition is not always the enemy.} Sometimes the Cantor’s cult is sanctioned – the local lord enjoys the music. The threat of Inquisition should ebb and flow with the story.
\end{itemize}

\noindent \textit{“Sing, and they will come – not just to listen, but to learn, to serve, to write down every note. And in that writing, the Patron grows fat. Feed carefully.”}

\clearpage

\chapter{Last Words from the Gray Wanderer}

You’ve read the mechanics. Now close the book and look at your character sheet. That Obligation number is not a debt to be minimised – it’s a \textbf{promise} that you matter enough to be bound.

When you stand before the Patron’s altar, lacking your Symbol, your Codex stolen, your voice hoarse, and you still reach for power using a broken coin and a desperate whisper – that is \textit{Fate's Edge}.

\vspace{1cm}
\noindent \textit{Keep your components dry, your voice warm, and your wit sharper than any Rite.}

\vspace{0.5cm}
\hfill — \textit{The Gray Wanderer}

\end{document}
